\chapter{High availability}
\label{ch:high_availability}

\section{Introduction to high availability}
\label{sec:ha_intro}

A venue is highly available when the failure of one instance does not stop
members trading. This chapter states what this venue must do when a component
fails, and what a member can observe while it happens.

The matching engine must reach the same result from the same sequence of orders.
An instance that takes over must therefore hold the orders its predecessor held,
and hold them in the same order. An order the venue has taken must survive the
failure of the instance that took it. Members are waiting on their orders while
a failure is handled, so the venue must resume quickly.

Recovery from the death of a process on a working machine is faster than
recovery from the loss of the machine itself. This chapter specifies the two
separately.

\subsection{Architectural patterns}
\label{subsec:ha_architectural_patterns}

Two patterns are in common use for fault tolerance, and this chapter uses both
names.

Under \textbf{active-passive}, one instance does all the work and a second
instance holds a copy of its state. The second does none of the venue's work
while the first is running. When the first fails, the second takes over.

Under \textbf{active-active}, several instances process transactions at the same
time. They agree on the order of those transactions through a consensus protocol
such as Raft or Paxos, or through synchronous state machine replication.

This venue is active-passive. \Cref{subsec:ha_approach} says which components
that applies to.

\subsection{What a failover must preserve}
\label{subsec:ha_failover_mechanics}

A failover affects three kinds of state. This chapter specifies each of them
separately.

\begin{enumerate}
    \item \textbf{Session state.} A member's session must survive the failure of
      the component holding it, together with its sequence numbering and the
      terms the member was admitted under.
    \item \textbf{Order book state.} The instance that takes over must hold the
      same resting orders as the instance that failed. If it holds fewer, orders
      that should have been filled are not filled. If it holds more, an order can
      be filled twice.
    \item \textbf{Gateway state.} A gateway translates a member's protocol and
      applies risk checks to what the member sends. When a gateway fails, the
      member must be able to establish whether an order it had sent was taken by
      the venue.
\end{enumerate}

\subsection{Two properties, and why one does not imply the other}
\label{subsec:ha_two_properties}

High availability is usually discussed as one property and is two. They fail
independently, and a venue can have the first in full while having none of the
second.

\textbf{Continuity} is the venue continuing to work. Something detects that an
instance has stopped, another instance takes over, and members go on placing
orders. Almost everything written about failover is about this: how it is
detected, how long it takes, how the decision is reached.

\textbf{Custody} is the venue keeping what it was given. An order accepted before
the failure is still there afterwards, on the same terms; a session's numbering
continues; a report the venue had not yet sent can still be sent.

\textbf{Continuity is easy to observe and custody is not}, which is what makes
the pair dangerous. A venue that has lost everything it was holding looks
entirely healthy from outside: it accepts orders, answers them promptly, moves
leadership when an instance dies, and reports no error. The evidence of lost
custody is an absence --- an order that stops being mentioned --- and an absence
is not something an observer notices, nor something a test detects unless it was
written to look for that particular one.

The two are also easy to conflate when reasoning. A failover that completes in
three seconds with orders flowing throughout is a success by every measure of
continuity, and says nothing whatever about whether the promoted instance holds
what its predecessor held.

\textbf{This is why the four questions are shaped as they are.} Only the last of
them --- what the member is told --- concerns what an observer would see. The
first three are about state: what a component holds, where the copy that outlives
it lives, and how a restarted instance gets back to current. Continuity follows
from a component running again; custody does not, and has to be specified.

\subsection{The pattern this venue uses}
\label{subsec:ha_approach}

This library is intended for low latency environments rather than ultra-low ones.

\textbf{This system is active-passive} for every component that must be
singular: the sequencer, the matching engine and the arbiter. Each of those runs
as a pair, one instance the leader and the other the follower. The follower does
none of the venue's work while it holds that role, and nothing it holds reaches a
member until it becomes the leader.

\begin{itemize}[nosep]
  \item \textbf{The leader does the work.} The sequencer assigns each order a
        sequence number, appends it to the write-ahead log and forwards it to the
        matching engine. The matching engine accepts the order and issues the
        execution report for it.
  \item \textbf{The leader replicates continuously.} The sequencer streams every
        log record to its follower. The follower writes each record to a log of
        its own and acknowledges it. The matching engine's follower keeps a
        replica of the orders it holds open, taken from the same stream.
  \item \textbf{The follower acts on nothing.} It forwards no order, issues no
        execution report and answers no gateway. While it holds that role it
        processes no traffic at all.
  \item \textbf{On failure the follower is promoted.} The arbiter tells it that
        it now leads. It brings itself up to date from whatever it has replicated
        but not applied, and then begins acting.
\end{itemize}

That is the pattern \cref{subsec:ha_architectural_patterns} describes as
active-passive. At no moment do two instances process the same traffic, and no
transaction is ever handled by more than one instance at a time.

An arbitration service and a witness decide which instance may act. Together they
stop both instances believing they lead at the same time.

Components that need not be singular run differently. The authentication
services run active-active, and any instance answers any caller. Several
gateways run at once, each holding a different set of sessions, which divides
the work rather than replicating it. \Cref{sec:ha:disabled:scope} returns to
that distinction, because it decides what turning high availability off
removes.

\section{Preliminaries}
\label{sec:ha:preliminaries}

A requirement is a statement of behaviour the venue must exhibit. It carries a
permanent identifier of the form \texttt{R-0001}, allocated in sequence and
never reused, so that a citation from a test, a defect record or a commit
message resolves for the life of the project. The identifier does not encode
position: a cross-reference renders as the identifier followed by the section it
is stated in, computed when the document is built.

A requirement describes observable behaviour --- what the venue does, and where
a member is involved, what that member is able to observe or do. Something that
holds when nothing has failed is an assumption rather than a requirement, and is
stated as one. One requirement is stated for each outcome a member can
distinguish, rather than for each mechanism that produces one.

\textbf{A requirement is stated of the venue} unless it is genuinely about one
component, and stating it of the venue would change what it means.
Naming the component that must do something is a decision about how the venue is
built, and this chapter does not make those. Where a requirement applies to every
component separately it says \emph{a component}, which is a different claim from
\emph{the venue} and is used where that difference is meant.

\textbf{This chapter says what the venue does, and not how it does it.} Where a
requirement constrains what may happen while a component waits for a disk, or
what a component must still be doing while it waits, it says only that.
\Cref{ch:library} decides which data structures satisfy such a requirement, and
this chapter does not name them.

\textbf{This chapter was written after the components it describes.} The venue
was built first, and its behaviour under failure was settled case by case as
each case arose. That produced an implementation whose parts had never been
reconciled with one another. This chapter states what those components must do,
arrived at by asking the same questions of each of them. Where a component does
something different today, a gap says so.

A \textbf{gap} records something the venue does not do that this chapter says it
must. Where the difference has already been recorded as a defect, the gap names
it. One is cited often enough to be worth stating once here:

\begin{itemize}[nosep]
  \item \textbf{BUG-0048} --- nothing bounds how far back the venue's record of
        what it did is kept. The venue discards none of it, so the record grows
        for as long as the venue runs, and how far back it reaches is whatever
        the disk has allowed rather than a period the venue states.
\end{itemize}

Every requirement names the scenarios that verify it, or states that none does
yet. That note says why no scenario exists; where the reason is that the venue
does not do the thing at all, a gap says so separately. The two are kept apart
deliberately: a requirement nothing tests and a requirement the venue does not
meet are different conditions, and only the second is a defect. \textbf{A
requirement that nothing verifies is still complete.} What is missing is a test
of the venue. The count of unverified requirements therefore measures the venue
rather than this document, and it falls only as tests are written.

\bigskip

Components are started by a launcher, which restarts what it starts and knows
nothing else. It takes no part in deciding which instance holds which role, and
no requirement in this chapter depends on it being present.

A component that dies repeatedly has not failed once; it is failing. Restarting
it again produces another failure, so the number of attempts is bounded, and
exhausting that bound is a terminal condition rather than a slower loop. This
condition is called flapping, and it is specified in
\cref{sec:ha:disabled:supervision}.

\textbf{An order is committed} when the sequencer has assigned it a sequence
number and appended it to the write-ahead log. That moment divides every
scenario in this chapter. Before it, the venue has not taken the order: it may
be lost by a failure, and the requirements say only that it must not take effect
afterwards and that the member can establish what happened. After it, the venue
has taken the order and must answer it, whatever dies subsequently. The
word is never used as a synonym for the \texttt{New} execution report that
reports an order's acceptance to a member.

The record it is appended to is a write-ahead log: every order is written to it
before anything acts on the order, and it is the venue's account of what it has
taken. Where two instances of the sequencer run, an order becomes durable in two
stages rather than one --- written by the instance that took it, and acknowledged
by the second --- and a member is not told about an order until both have
happened. \Cref{sec:ha:enabled:sequencer:durable} states that rule and what it
is for.

\textbf{The following hold when nothing has failed}, and are specified
elsewhere. They are recorded here because the requirements in this chapter rely
on them.

\begin{itemize}[nosep]
  \item A member's session is its comp id and protocol rather than its
        connection, and outlives the connection that carried it.
  \item Sequence numbering is reset on a stated schedule, and never as a
        consequence of a failure.
\end{itemize}

\section{High availability disabled}
\label{sec:ha:disabled}

The venue has a single switch that turns high availability off. With it off no
arbiter and no witness is started, and each component that would otherwise run
as a primary and secondary pair runs as one instance.

This configuration exists for development, and for no other purpose. A
development machine is one machine, and a second instance of every component
consumes processors and memory that the work being done has better uses for.
Many tests do not depend on high availability at all, and run faster and more
predictably without the components that provide it.

\textbf{It is not a supported production configuration and must never be
deployed as one.} With one instance of each singular component there is nothing
to fail over to. The only recovery available is to restart the process that died
and rebuild what it held.

\textbf{Most testing runs with high availability enabled.} The scenarios that
exercise the behaviour specified in this chapter run in that configuration,
because it is the configuration a production venue runs. This configuration is
tested for what is particular to it: that the components which exist only to
arbitrate are absent, and that the components which remain report their absence.

\textbf{Development is not restricted to this configuration.} High availability
can be switched on with every component running on one machine, and doing so
starts the arbiters, the witness and a second instance of each singular
component exactly as a production deployment would. A developer can therefore exercise the behaviour specified in
\cref{sec:ha:enabled} on one machine. Such a deployment cannot reproduce the loss
of a machine, because it has only one, so it exercises everything in that section
except \cref{sec:ha:enabled:machine_loss}.

\subsection{What the switch does and does not remove}
\label{sec:ha:disabled:scope}

Two kinds of multiplicity must not be confused.

\emph{Replicated peers} are components of which several instances run, all of
them live, none of them leading. The authentication services and the gateways
are replicated peers: a caller reaching any instance gets the same answer, and
losing one costs capacity rather than correctness.

\emph{Singular components} are those of which exactly one instance may act at a
time. The sequencer, the matching engine and the publisher are singular, and
that constraint is what the arbiter, the leader and follower roles, and the
leadership generation exist to enforce.

Turning high availability off removes the second instance of each singular
component, and the components whose only purpose is to arbitrate between them.
It says nothing about the replicated peers, which run in every configuration.

\subsection{What this configuration may and may not do}
\label{sec:ha:disabled:obligation}

This configuration is used for development and for running tests. Two
requirements follow from that, and they constrain the configuration in opposite
directions.

\textbf{This configuration may lose what a production venue must not.} It runs
one instance of each singular component, so the death of that instance can lose
an order the venue has taken, and can lose a session's numbering with it. That
follows from running one instance of everything, and it is accepted here because
nothing of value is being traded.

\textbf{It must not behave differently from production in ways a test cannot
see.} A test run against this configuration is evidence about the venue only if
the venue does the same things here that it does in production, and fails for the
same reasons. Where it does less, it must fail visibly rather than appear to
work.

The failure this rule prevents is a venue that accepts a member's orders,
acknowledges every one of them, forwards none of them anywhere, and logs no
error. A test that sent orders and received replies would pass against such a
venue, and the result would mean nothing.

\subsection{Recovery in general}
\label{sec:ha:disabled:recovery}

Two requirements hold for every component in this section, and are stated once
here rather than repeated in each. Both concern the window between a component
starting and its being ready, which is short, is entered on every restart, and
is where a second failure does the most damage.

Every requirement here assumes that what survived a failure can be read. A log
truncated part-way through a record, a checkpoint that fails its checksum, or a
file that is simply absent leaves a component with nothing to recover from --- and
in this configuration with nowhere else to turn, because there is no second
instance holding a clean copy.

\textbf{This configuration is weak there, and deliberately so.} With high
availability enabled the instance whose copy is intact takes over, which is the
answer to damaged state and the reason the enabled configuration exists. Nothing
in the disabled configuration substitutes for it. That is a reason to run this
configuration only where losing what it holds is acceptable.

\begin{gap}{}
  Damaged durable state is not detected as such. A component reads what it can
  and proceeds, rather than establishing that what it read was whole, so a
  truncated log or a partial checkpoint produces a component that holds the wrong
  orders and reports itself as current. A member is then answered from that
  state: an order it placed may be unknown, and one it cancelled may be live.
\end{gap}

\begin{requirement}{R-0121}{A component takes the cost of first touching its memory before it reports itself ready}
  Where a component maps memory it will use while trading, it shall touch all of
  that memory before reporting itself ready, so that the delay of first use falls
  before it is trading rather than during.
  \rationale{Mapping a file reserves addresses and reads nothing. The first time
  each page is used the processor interrupts and the kernel fetches it --- a small
  delay, once per page, and for a region of a few hundred megabytes there are tens
  of thousands of them. Left alone they fall on whatever code touches each page
  first, which is the order path.

  Doing it before reporting ready means a component that has said it can trade
  has already paid them. Doing it afterwards moves the cost onto the first
  members to send anything, which is the worst moment to take it and the hardest
  to attribute afterwards.}
  \uncovered{no scenario measures a component's first orders against its later ones}
\end{requirement}

\begin{requirement}{R-0122}{A report a member receives twice is marked as a repeat}
  Where a recovery repeats work whose execution reports had already been sent,
  the reports sent again shall be marked so that the member can tell a repeat
  from a second event.
  \rationale{A component that dies partway through handling an order cannot know
  whether the report for it was sent. Sending it again is the right choice: a
  member that never receives one cannot recover from that alone, and a member
  that receives one twice can, provided it can tell.

  Unmarked, a repeated acceptance is indistinguishable from a second order
  accepted --- so the venue would be reporting something that did not happen, in
  the course of correcting something that did.}
  \covers{ha_test:53}
\end{requirement}

\begin{gap}{}
  Scenario 53 covers the repeat a catch-up produces, which is the common case and the one a
  member meets: an engine applies the venue's record and marks every report it sends, because
  it cannot tell a record an earlier engine reported from one no engine ever saw.

  It does not cover the narrower case the requirement was written for --- a component that dies
  between sending a report and recording that it did, so that the same work is repeated on the
  strength of a record that is one step behind. Nothing kills a component in that window, and
  producing it deliberately is harder than the other scenarios: the window is a few
  instructions wide.
\end{gap}

\begin{requirement}{R-0031}{A component does not act on partially recovered state}
  A component shall not begin work until its recovery is complete. No sequence
  number shall be assigned before the last one issued is known, and no order
  shall be accepted before the set of open orders is current.
  \rationale{Acting on partial state produces results that look ordinary and are
  wrong. Of the ways a venue can fail, that is the hardest to detect and the
  hardest to unwind afterwards.}
  \uncovered{no scenario interrupts a recovery}
\end{requirement}

\begin{requirement}{R-0032}{An interrupted recovery can be repeated without applying anything twice}
  Where a component fails while recovering, the recovery shall be able to run
  again from the same starting point and reach the same result.
  \rationale{Without this a failure during recovery corrupts rather than delays,
  and a component that fails repeatedly leaves worse state on each attempt. It
  is also what makes a bound on the number of attempts meaningful: the bound is
  only safe if the attempts beneath it were harmless.}
  \uncovered{no scenario interrupts a recovery}
\end{requirement}

\subsection{The authentication service}
\label{sec:ha:disabled:auth}

Several authentication services run at once, and none of them leads. A member
never contacts one directly. A gateway contacts an authentication service on the
member's behalf at logon, and reaches whichever instance answers.

Losing one instance reduces the capacity available and does not affect
correctness. Two instances holding different answers to the same question does
affect correctness, and most of this subsection concerns that condition.

\subsubsection{What state it holds}
\label{sec:ha:disabled:auth:state}

An instance holds the credential set: which comp ids may trade, the material each
is authenticated against, and the terms each was admitted under. This set is a
copy rather than the original. The database holds the original, and the
administration service is the only component that writes to it. That service
changes credentials and distributes each change to every authentication service,
and \cref{sec:ha:disabled:admin} describes it.

\subsubsection{Where the copy that outlives it lives}
\label{sec:ha:disabled:auth:durable}

The database holds the copy, and an instance reaches it in two steps rather than
one. An instance loads an export of the database when it starts, and the
administration service afterwards sends it each change as it makes it. The
instance holds no connection to the database, so it can neither check an answer
against the original nor establish how far behind it has fallen.

\textbf{Only the administration service can establish whether an instance is
current}, because only that service holds a record of what it sent and to which
instance.

\begin{gap}{}
  A change is reported as having succeeded when at least one instance applied
  it. A partial fan-out is therefore indistinguishable from a complete one, and
  nobody learns that an instance was missed.
\end{gap}

\begin{gap}{}
  Nothing brings a running instance up to date. Having missed a change, it stays
  wrong until it is restarted --- and restarting reloads the export, which may
  itself predate the change, so restarting is not a reliable remedy either.
\end{gap}

\subsubsection{How a restarted instance becomes current}
\label{sec:ha:disabled:auth:recovery}

The instance reloads the export and resumes answering. It has one such operation
rather than two, as the gateway in \cref{sec:ha:disabled:gateway} does, because
an authentication service holds no role and is never promoted. A restarted
instance and a newly started one therefore do the same thing.

The instance does not ask for the changes it missed, and no component offers them
to it.

\subsubsection{What the member is told}
\label{sec:ha:disabled:auth:member}

An instance can be wrong in two ways, and the two are not equally serious, which
is why the requirements below treat them differently. An instance that has missed
an addition refuses a member that should be admitted. The member sees that
refusal at once, and the venue can recover from it. An instance that has missed a
withdrawal admits a member that should be refused. Nothing reports that
admission, and it persists until somebody corrects the instance.

\begin{requirement}{R-0042}{A credential change is applied by every instance, or is reported as incomplete}
  A change shall not be reported as having succeeded unless every instance
  applied it. Where an instance did not, the change shall be reported as
  incomplete and shall name the instance.
  \rationale{Reporting success when one instance of several applied the change leaves a
  partial distribution looking the same as a complete one, so nobody can act on it.
  Every remedy below depends on the venue knowing that the condition exists.}
  \uncovered{no scenario makes a credential change while an instance is down}
\end{requirement}

\begin{requirement}{R-0043}{An instance that has missed a change is brought current without being restarted}
  Where an instance has missed a change, it shall be brought up to date while
  running.
  \rationale{Restarting an instance is not a remedy, because the instance reloads an export of
  the database that may itself be older than the change it missed. Nothing else
  brings an instance up to date, so without this requirement an instance stays
  stale indefinitely.}
  \uncovered{no scenario makes an instance miss a change, and nothing would bring it current}
\end{requirement}

\begin{requirement}{R-0044}{A member is admitted by every instance or by none}
  A member's logon shall succeed or be refused alike whichever instance answers
  it.
  \rationale{Intermittent admission is worse for a member than consistent refusal. The logon
  succeeds, then it fails, and nothing distinguishes the two attempts, so the
  member can neither report the fault clearly nor reproduce it.}
  \uncovered{no scenario authenticates the same member against both instances}
\end{requirement}

\begin{requirement}{R-0045}{An instance can say how current its credential set is}
  An instance shall be able to report the point up to which it has applied
  changes.
  \rationale{Where an instance cannot report this, no component can establish it. The instance
  has no view of the database, and the administration service cannot distinguish an
  instance that missed a change from one that never received it.}
  \uncovered{no scenario asks an instance how current its credentials are, and nothing answers}
\end{requirement}

\begin{requirement}{R-0046}{A change made while an instance is unavailable is applied when it returns}
  A credential change shall be accepted while an instance is unavailable, that
  instance shall be recorded as not current, and the change shall be applied to
  it when it returns.
  \rationale{Refusing the change instead would stop the venue withdrawing access during an
  outage, which is when withdrawing access may matter most. The cost is an instance
  that is behind for a bounded period. That cost is acceptable because
  \reqref{R-0042} makes the venue aware of the condition.}
  \uncovered{no scenario makes a credential change while an instance is down}
\end{requirement}

\subsection{The administration service}
\label{sec:ha:disabled:admin}

The administration service is the only component permitted to change credential
state. Members never contact it. Authentication services receive what it sends
them.

\textbf{One instance runs in both configurations, by decision.} The service is
not on the path an order takes, so its absence stops no member trading or logging
on. Credentials change rarely, so the venue tolerates its absence for as long as
a restart takes. A second instance would need arbitration of its own to remain
the single writer, which is a large mechanism for a component that is idle for
most of a trading day.

That decision carries one condition: everything the single instance must retain
has to survive the instance.

\subsubsection{What state it holds}
\label{sec:ha:disabled:admin:state}

\textbf{The credential state.} The service writes this state, and the database
holds the original.

\textbf{Which authentication services have applied which change.} No other
component holds this. An authentication service holds no record of a change it
never received, and the database records what the credentials are rather than
which instance has been told.

\subsubsection{Where the copy that outlives it lives}
\label{sec:ha:disabled:admin:durable}

The database holds the credential state. Nothing holds the record of which
instances have applied which change.

\begin{gap}{}
  Which instances have applied a change is held only in the running process. An
  administration service that restarts between making a change and distributing
  it to every instance cannot afterwards establish which instances were missed.
\end{gap}

\subsubsection{How a restarted instance becomes current}
\label{sec:ha:disabled:admin:recovery}

The service reads the credential state from the database, and that state is
complete. It cannot recover the record of what it had distributed, because
nothing wrote that record down.

The service has one such operation rather than two, as the gateway in
\cref{sec:ha:disabled:gateway} and the authentication service in
\cref{sec:ha:disabled:auth} do. An administration service is never promoted,
because the venue never runs a second one to promote.

\subsubsection{What the member is told}
\label{sec:ha:disabled:admin:member}

Nothing directly. An authentication service answers a member's logon, and the
requirements in this subsection exist so that the answer does not depend on an
administration service running at the time.

\begin{requirement}{R-0096}{A credential change is recorded before it is distributed}
  A change shall be written to the source of truth before it is sent to any
  instance, so that a failure between the two loses the distribution rather than
  the change.
  \rationale{The next instance to reload from the database undoes a change that was
  distributed but not recorded, and nothing reports that. A change that was
  recorded but not distributed appears as a difference between instances, and the
  venue can correct it.}
  \uncovered{no scenario interrupts a credential change}
\end{requirement}

\begin{requirement}{R-0097}{The record of which instances are behind outlives the service that made it}
  Where an instance did not receive a change, that shall be recorded durably, and
  shall be acted on by whichever administration service is running when the
  instance next becomes reachable.
  \rationale{\reqref{R-0046} accepts a change while an instance is unavailable, on the basis
  that the venue applies it later. Where only a running process holds the
  obligation to apply that change, the obligation is lost whenever the
  administration service restarts before the instance returns. The credential state
  is then permanently inconsistent, and nothing records that it is.}
  \uncovered{no scenario restarts the administration service}
\end{requirement}

\begin{requirement}{R-0112}{The venue trades while the store it depends on is unreachable}
  No component shall require a connection to the durable store in order to admit
  a member, accept an order, or answer one. What it needs from the store shall be
  read when it starts and supplied to it thereafter.
  \rationale{Where the venue needs the store to be reachable in order to work, every outage of
  the store becomes a trading outage, and the venue takes on a failure whose
  likelihood has nothing to do with anything the venue does. Admission already
  works this way: an authentication service reads an export when it starts and is
  sent each change afterwards. Orders that outlive the day are to work the same
  way.}
  \uncovered{no scenario runs the venue with the store unreachable}
\end{requirement}

\begin{requirement}{R-0098}{Members trade and log on while no administration service is running}
  The absence of an administration service shall not prevent a member logging on,
  placing orders, or being answered.
  \rationale{The administration service writes credential state, and answers nothing when a
  member is admitted. Making a member's logon depend on it would place a component
  that is idle for most of a trading day on the path of every session that starts.}
  \uncovered{no scenario runs the venue without an administration service}
\end{requirement}

\subsection{The gateway holding the session}
\label{sec:ha:disabled:gateway}

The gateway carrying a member's session dies and is restarted. The member's
connection ends and it reconnects.

\subsubsection{What state it holds}
\label{sec:ha:disabled:gateway:state}

The gateway holds the state of the session: how far the session's numbering has
reached in each direction, and which of the outbound numbers carried execution
reports rather than session-level traffic. This is the only state that matters
when the gateway dies.

While it is running, the gateway also holds the member's identity and the
instructions the member gave when it was admitted --- whether its resting orders
should be cancelled if its connection ends, and how long the venue should wait
before doing so --- any authentication exchange part-completed, any
resend being served, and any reports held back while the numbering is being
established. None of these needs to survive the gateway. The member's identity is
established again at the next logon, and the venue abandons the rest, which the
member asks for again.

The gateway holds no order state. The matching engine holds which orders are
open, and the sequencer holds the order they were taken in. Where the gateway
must act on a member's behalf, such as cancelling what the member has resting
when its connection ends, it reports the disconnection and those instructions to
the component that owns the orders. It keeps no view of those orders to act on
itself. The instructions therefore belong to the session rather than to the
gateway.

Those instructions reach the gateway when the member is admitted, and a component
needs them later, possibly a different instance, and possibly when the member is
no longer connected to supply them again. The requirement \reqref{R-0103} states
what must be true wherever they are held, and this chapter does not say where
that is. \texttt{docs/availability/session\_binding.md} makes that choice and
gives the reason.

Two cases have to be told apart, and only one of them is answered by holding the
instructions somewhere a gateway can ask. A gateway that restarts is asked for
nothing until the member binds again, and binding is what fetches them. A gateway
that dies holding live sessions is asked nothing at all, and a member that never
reconnects binds nowhere --- so the instruction that was to fire at that moment
must be carried out by whatever owns the orders, which is \reqref{R-0113}.

\begin{gap}{BUG-0091}
  The instructions are held only by the gateway that was given them. A gateway
  that dies takes with it the instructions for every session it held, and there
  is nowhere else to read them from.
\end{gap}

\subsubsection{Where the copy that outlives it lives}
\label{sec:ha:disabled:gateway:durable}

The sequencer holds the numbering. A gateway reports a session's position as it
advances, and the sequencer records it, because the sequencer is the only
component that every gateway instance talks to. The sequencer's write-ahead log
holds the execution reports, each stamped with the session it belongs to.

An order that has not been committed has no copy anywhere in the venue. Between
the member sending it and the sequencer committing it the only record inside the
venue is the gateway's, and when the gateway is what died, nothing can recover
it.

This follows from a deliberate decision. A durable record kept by the gateway
would add a write to the hot path, in front of the write-ahead log that already
exists, and would double the cost of every order to protect a window of
microseconds. The member that sent the order also holds a copy of it outside the
venue. The venue therefore does not attempt to recover an uncommitted order. It
guarantees that such an order never takes effect afterwards, that the member can
establish this, and that the member can send the order again.

\begin{gap}{BUG-0090}
  The gateway keeps its own view of what each session has resting, and acts on
  that view when a connection ends. The view is built only as execution reports
  arrive, so a restarted gateway holds nothing for orders placed before the
  restart, and a member that asked for cancel-on-disconnect is no longer
  protected by it, is not told, and cannot tell.
\end{gap}

\begin{requirement}{R-0008}{Reports remain recoverable for the reconnection window}
  The venue shall retain every execution report for at least as long as a member
  may take to reconnect and ask for it, and that period shall be stated.
  \rationale{Recovery that depends on a record still being present is only as good as the
  retention behind that record. A retention period that follows from how much
  storage is available, rather than from a stated figure, gives a member nothing to
  rely on.}
  \uncovered{no scenario measures how far back reports remain available}
\end{requirement}

\begin{gap}{BUG-0048, BUG-0046}
  There is no stated retention period. The write-ahead log discards nothing, so a
  report stays available for as long as the venue runs and the disk holds out. A
  member that reconnects is therefore served by what the venue happens still to
  hold, rather than by a period the venue has undertaken to keep.
\end{gap}

\subsubsection{How a restarted instance becomes current}
\label{sec:ha:disabled:gateway:recovery}

The member reconnects, the session binds to the restarted gateway, and the
sequencer supplies the numbering it holds. What the member receives next depends
on how far each of its orders had got before the gateway died.

The gateway has one such operation where other components have two, because a
gateway is never promoted and holds no role. A member returning to the instance
that restarted, and a member arriving at a different instance, are the same event
from the venue's side: a session binds to an instance that does not hold its
state, and that instance asks the sequencer for it. Nothing in that sequence
depends on how many gateways are running, so gateways need no arbiter and running
several of them costs little.

\textbf{An order the gateway had received but not forwarded} was never
committed.

\begin{requirement}{R-0001}{An order that was not forwarded never takes effect}
  The order shall have no effect on the venue. No execution report shall be
  produced for it, and the venue shall hold no order open on its behalf.
  \rationale{The venue had not taken the order. Acting on it later would execute an
  instruction that the member has been given no reason to believe is live.}
  \uncovered{no scenario restarts a gateway}
\end{requirement}

\textbf{An order the sequencer had committed} is the venue's responsibility from
that moment.

\begin{requirement}{R-0004}{A committed order is processed}
  The order shall be processed by the matching engine, whether or not the
  gateway that received it is running when the engine reaches it.
  \rationale{Committing the order is the venue taking it. What happens after that point must
  not depend on which component carried the order.}
  \uncovered{scenario 23 covers a gateway that is not restarted}
\end{requirement}

\textbf{A report that had been given an outbound sequence number} but did not
reach the member leaves a gap in the member's numbering.

\begin{requirement}{R-0007}{A numbered report that was not received is recoverable}
  The member shall detect the gap in its inbound numbering and ask for the
  missing range, and the venue shall send again the execution reports that
  occupied those numbers, filling the numbers that carried session traffic
  rather than reports.
  \rationale{This is the FIX resend, used as the protocol intends. The venue supplies the
  reports it can reproduce, and accounts for the numbers it cannot.}
  \covers{ha_test:22}
\end{requirement}

\subsubsection{What the member is told}
\label{sec:ha:disabled:gateway:member}

\begin{requirement}{R-0103}{The instructions a member gave at logon outlive the component that received them}
  The instructions a member gave about what should happen to its orders when its
  connection ends shall be available to whichever component must act on them,
  whether or not the component that first received them is still running.
  \rationale{Those instructions arrive once, when the session is admitted, and a component
  acts on them later, possibly a different instance, and possibly when the member
  is no longer present to supply them again. Instructions held only by the
  component that received them are lost in exactly the failure they guard against.}
  \uncovered{no scenario restarts a gateway holding a session's instructions}
\end{requirement}

\begin{requirement}{R-0113}{A member's instruction is carried out when the component holding it is the one that failed}
  Where a member has asked for its orders to be cancelled if its connection
  ends, and the connection ends because the component holding that session
  failed, the instruction shall still be carried out.
  \rationale{A component acts on the instruction by reporting the disconnection to whatever
  holds the orders, and a component that has died reports nothing. The death of a
  gateway ends every session it holds, and is therefore the one failure after which
  nothing acts on what those sessions asked for. The member cannot detect this,
  because from the member's side the connection ended in the ordinary way.}
  \uncovered{no scenario ends a session by killing the component holding it}
\end{requirement}

\begin{requirement}{R-0002}{A member can establish that an order was not taken}
  On reconnecting, a member shall be able to determine that the venue holds no
  order for a \texttt{ClOrdID} it had sent.
  \rationale{The member cannot distinguish this case from a lost reply. Every
  venue provides an order status enquiry for exactly this reason.}
  \uncovered{no scenario asks the venue what it holds for a member, because nothing answers}
\end{requirement}

\begin{gap}{BUG-0089}
  The venue answers no order status enquiry of any kind, so a member has no way
  to ask what is held for it. Recovery rests entirely on what the venue chooses
  to send.
\end{gap}

\begin{requirement}{R-0003}{An order that was not taken may be resubmitted unchanged}
  A member may send the order again under the same \texttt{ClOrdID}, and the
  venue shall process it as a new order.
  \rationale{Without this requirement, the member must choose between abandoning an order and
  risking a duplicate. Resubmitting under the original identifier removes that
  choice.}
  \uncovered{no scenario restarts a gateway}
\end{requirement}

\begin{requirement}{R-0006}{A resubmitted committed order is refused as a duplicate}
  Where a member sends again an order the venue had already committed, under the
  same \texttt{ClOrdID}, the venue shall refuse it as a duplicate and shall not
  create a second order.
  \rationale{This requirement is what makes \reqref{R-0003} safe. A member cannot tell the two
  cases apart, and relies on the venue to distinguish them.}
  \uncovered{no scenario restarts a gateway}
\end{requirement}

\begin{requirement}{R-0010}{A duplicate is refused whatever has become of the first order}
  The refusal required by \reqref{R-0006} shall not depend on the order having
  reached the matching engine, nor on the engine that received it still running.
  \rationale{The sequencer is the component that commits an order, so only a record of what
  has been committed can answer this question at the moment it is asked. The
  matching engine does not hold an order that has been committed and not yet
  applied, and an engine that has restarted holds no record of an earlier
  submission. In either case a member following \reqref{R-0003} would be given a
  second live order.}
  \uncovered{no scenario restarts a gateway}
\end{requirement}

\begin{requirement}{R-0119}{An order identifier used today is refused for the rest of the day}
  Where a member sends an order under an identifier it has already used on that
  session that trading day, the venue shall refuse it, whatever became of the
  first order and whatever has restarted since.
  \rationale{Two live orders under one identifier make a cancel request ambiguous, and a check
  against the orders currently open would catch that much. The reason to go further,
  and refuse an identifier whose first order has finished, concerns the record the
  venue publishes. Consumers of that record identify an order by this value, and a
  consumer such as a member's back office, a reconciliation process or an auditor
  cannot ask the venue anything. Two different orders under one identifier in that
  record cannot be read correctly.

  This requirement also makes \reqref{R-0003} safe. This document requires
  resubmission under the original identifier as a recovery path, for a member that
  does not know whether the venue took the order. The check that distinguishes a
  resubmission from a new order is therefore something a member following this
  specification depends on. If that check does not survive a restart, a member
  following the specification is given a second live order at the moment it is
  least able to notice.

  The trading day bounds the set of identifiers. The venue adds identifiers through
  a trading day and discards the set at the end of it, which keeps the set finite.}
  \uncovered{no scenario reuses an identifier whose order has been cancelled}
\end{requirement}

\begin{gap}{}
  An identifier is checked against the orders the venue currently holds open, and
  a cancelled order is removed from that set --- so an identifier whose order has
  been cancelled is accepted again, and the activity the venue publishes can carry
  the same identifier for two different orders on the same day.

  Nothing durable holds the identifiers used during a day either, so even the
  check that exists does not survive a restart of the component that makes it.
\end{gap}

\begin{requirement}{R-0005}{A report produced for an unbound session is delivered on reconnection}
  An execution report produced while the session had no connection shall be
  delivered when that session next binds to a gateway, without the member having
  to ask for it.
  \rationale{The venue allocated no outbound sequence number for such a report, so it leaves
  no gap in the member's numbering. The member has nothing to notice and nothing to
  request, and a resend cannot reach it.}
  \uncovered{no scenario restarts a gateway}
\end{requirement}

\begin{gap}{BUG-0088}
  Recovery is driven entirely by the member noticing a gap and asking. A report
  produced while the session had no connection was never numbered, so it leaves
  no gap, and nothing pushes it. The report is written to the log stamped with
  its session, so what is missing is the delivery rather than the record.
\end{gap}

\begin{requirement}{R-0009}{A gap that cannot be filled is reported as such}
  Where the venue cannot supply a report the member has asked for, it shall say
  so rather than fill the number silently.
  \rationale{A gap filled without explanation leaves the member believing its numbering is
  complete when it has lost a report it will never see.}
  \uncovered{no scenario restarts a gateway}
\end{requirement}

\subsection{The sequencer}
\label{sec:ha:disabled:sequencer}

The sequencer dies and is restarted. This breaks no member's connection, because
gateways hold the sessions and the gateways stay up. What stops is the venue's
ability to take an order.

\subsubsection{What state it holds}
\label{sec:ha:disabled:sequencer:state}

\textbf{The ordering.} The sequencer holds the number it will assign next, and
the write-ahead log that records what it has assigned. This log is the venue's
record of what it has taken. The rest of the state in this subsection depends on
it.

\textbf{The leadership generation.} The sequencer holds the generation its work
is stamped with, so that a receiver can distinguish a message from a superseded
instance from a message from the current one.

\textbf{Each session's numbering.} The sequencer holds how far every session has
reached in each direction, and which of its outbound numbers carried execution
reports. It holds this on the gateways' behalf rather than for its own use, as
\cref{sec:ha:disabled:gateway:durable} describes, because it is the only
component that every gateway instance talks to.

\textbf{Where an execution report should go.} The sequencer holds the route from
a committed order back to the session that placed it.

\textbf{What it has deferred.} The sequencer holds how many orders it has taken
that no matching engine has yet received, and how long they have been waiting.

\subsubsection{Where the copy that outlives it lives}
\label{sec:ha:disabled:sequencer:durable}

\textbf{The ordering and the generation survive the process.} The log is on
disk, and the sequencer takes the next number from the log's last record. The
generation is in a small file beside the log. The sequencer writes that file by
writing a temporary file and renaming it over the real one, so a reader sees
either the old value or the new one, whatever moment the process died at.

\textbf{The other three have no durable copy.} The running process holds them
and nothing else does. A restart treats each of the three differently.

The sequencer \emph{relearns} session numbering rather than recovering it.
Gateways report every session's position on a timer, so a restarted sequencer is
told again, within one reporting interval, about every session bound to a gateway
at that moment. No gateway reports a session that has no live connection, and
that is the session about to reconnect and ask for its numbering.

\begin{gap}{}
  A session that was not connected when the sequencer restarted is forgotten
  entirely. Nothing on disk holds its numbering and no gateway is reporting it.
\end{gap}

The sequencer does not rebuild the route from an order to a session, and that is
deliberate. Rebuilding it by replaying the log would fill it with entries that
nothing removes, and the log holds a trading day of them.

\begin{gap}{}
  An execution report arriving for an order committed before the restart cannot
  be placed, because the route it needed was not kept and is not rebuilt.
\end{gap}

What the sequencer has deferred is a count and a time. The orders themselves are
in the log, and a restarted sequencer holds no record that any of them is still
to be answered. It does not need one. A matching engine presents the position it
has reached, and the sequencer serves everything after that position from the
log, so which records an engine receives is determined by the position the engine
presents rather than by anything the sequencer remembers.

\subsubsection{How a restarted instance becomes current}
\label{sec:ha:disabled:sequencer:recovery}

The sequencer reopens the log, takes the next number from the last record it
finds, and reads the generation from the file beside the log. Other components
supply what it cannot recover: gateways report the sessions they hold, and a
matching engine connects.

What is outstanding at the moment the sequencer dies falls into two cases, and
the venue treats them differently.

\textbf{An order committed and not yet forwarded to the matching engine} belongs
to the venue, and \reqref{R-0013} covers it. The engine reports how far it has
reached, and the sequencer sends the records that follow.

\textbf{An execution report produced but not yet sent onward} is the second
case. The report survives, because the sequencer wrote it to the log stamped with
the session it belongs to. The route the sequencer was going to use does not
survive.

\begin{requirement}{R-0025}{A report can be addressed to its session after a restart}
  An execution report the venue has produced shall be deliverable to the session
  that placed the order, after a restart of any component between the two, and
  the session it belongs to shall be determinable from what survived that
  restart.
  \rationale{The record of the report already names its session. A route held only in a
  running process is lost at the moment it is needed, and no other component can
  supply it.}
  \uncovered{no scenario restarts a sequencer with reports undelivered}
\end{requirement}

\textbf{A promotion is a different operation, and it recovers more.} Gateways
report to both instances, so a follower has been receiving each session report
all along, and has been receiving every log record as it was written. A follower
therefore already holds what a restarted instance must be told, including the
numbering of sessions that are bound to no gateway, which are the sessions no
gateway will report.

The gateway in \cref{sec:ha:disabled:gateway} has one such operation, because a
gateway is never promoted. The sequencer has two. They recover different amounts,
and the one that recovers less is the one the venue uses when it has no second
instance to fail over to.

\subsubsection{What the member is told}
\label{sec:ha:disabled:sequencer:member}

\begin{requirement}{R-0026}{A logon completes or is refused within a bounded time}
  A member logging on shall be given either an established session or a refusal,
  within a stated time, whatever the venue has restarted meanwhile. It shall not
  be left authenticated and unestablished.
  \rationale{A session that is neither established nor refused leaves the member unable to
  trade and with nothing to act on. The member cannot proceed, and it has been
  given no reason to reconnect.}
  \uncovered{no scenario restarts a sequencer during a logon}
\end{requirement}

\begin{requirement}{R-0027}{An interrupted resend leaves the member able to act}
  Where the venue cannot finish supplying a range of messages it has begun, it
  shall say so. A member shall not be left with an unfilled gap and no
  indication that nothing further is coming.
  \rationale{From the member's side, waiting for records that will never arrive looks the same
  as waiting for records that are slow. Only the venue can distinguish the two
  cases, so only the venue can end the wait.}
  \uncovered{no scenario restarts a sequencer during a resend}
\end{requirement}

\begin{requirement}{R-0011}{A session's numbering survives a restart of whatever remembers it}
  A member reconnecting after the venue has restarted the component that holds
  its numbering shall be resumed where it left off, and not at the beginning.
  \rationale{A member whose numbering restarts sees a break it cannot reconcile, and the venue
  cannot tell that member which of its orders the venue still holds.}
  \uncovered{no scenario restarts a sequencer with a session unbound}
\end{requirement}

\begin{requirement}{R-0012}{A member is never resumed at a numbering the venue cannot vouch for}
  Where the venue has no record of a session's numbering, it shall say so rather
  than resume that session at the beginning.
  \rationale{A member whose own store has also restarted cannot distinguish a session resumed
  at the beginning from a session that has never traded. Both sides then agree on a
  numbering that neither has any basis for, neither side sees a gap, and the
  member's numbering is reset while its orders are still open. A member can recover
  from being refused. A member cannot recover from being agreed with wrongly.}
  \uncovered{no scenario restarts a sequencer with a session unbound}
\end{requirement}

\begin{gap}{}
  A member whose session the venue has forgotten is resumed at the beginning and
  is not told.
\end{gap}

\begin{requirement}{R-0013}{An order committed before a sequencer restart still receives its execution report}
  An order that the sequencer had committed before it died shall be processed,
  and an execution report shall be produced for it and be deliverable to the
  session that placed the order.
  \rationale{Committing the order is the venue taking it. Restarting the component that took
  the order does not return the order to the member.}
  \uncovered{no scenario restarts a sequencer with orders committed and unanswered}
\end{requirement}

\begin{requirement}{R-0014}{A sequence number is never issued twice}
  The venue shall not assign a sequence number it has assigned before, whatever
  has restarted since.
  \rationale{The sequence number is the venue's name for an order. The venue cannot answer two
  orders that share one number, and cannot distinguish the second from a
  retransmission of the first.}
  \uncovered{scenario 14 covers a restart with a peer to synchronise from}
\end{requirement}

\subsection{The matching engine}
\label{sec:ha:disabled:me}

The matching engine dies and is restarted. No member's connection is broken, and
the venue can still take orders. What stops is the processing of those orders.

\subsubsection{What state it holds}
\label{sec:ha:disabled:me:state}

\textbf{The set of open orders.} The engine holds which orders the venue has
accepted and has not yet cancelled or rejected, and the session each order
belongs to. The engine uses this set to recognise an order identifier that is
already in use, and to answer a cancel request. This set is the only copy of it.
The sequencer holds a record of which orders were taken and in what order, which
is a different thing, and converting that record into this set is expensive.

\textbf{The point the set has reached.} The engine holds the sequence number up
to which the set reflects what the sequencer has committed. Without that number
the engine cannot establish which committed orders are missing from the set.

\subsubsection{Where the copy that outlives it lives}
\label{sec:ha:disabled:me:durable}

In a memory-mapped region of fixed-size records, one for each open order, held
beside the engine and separate from the structures it works with. An order is
written to it as it is accepted and taken out of it as it is cancelled, and it is
read back when the engine starts. See
\texttt{docs/durability/open\_order\_checkpoint.md}.

The region survives the death of the process. It does not survive the death of
the machine, and this configuration is specified to give no more than that.
Where high availability is enabled, the answer to a lost machine is the other
machine.

\textbf{Reading the region is what a restart costs.} Scenario 50 recovers a
thousand open orders from a region of 248 MB in 56 milliseconds, where the pages
the dead process wrote are still held in memory. The same walk takes 457
milliseconds where those pages have to be read from the disk. Both figures follow
from the size of the region rather than from the number of orders open in it.

\textbf{The region holds state, and the sequencer's log holds changes.}
Rebuilding the set of open orders from the log instead means replaying the log
from the start of the trading day. The size of the region follows from the number
of orders open at one moment. The cost of a replay follows from the number of
orders the day has seen, which is far larger and grows all day, so a replay at a
full day's volume takes tens of seconds where reading the region takes
milliseconds.

The venue does replay the log in one case, which \reqref{R-0123} states: where
the region cannot be used at all. It does not resume trading while it does so. It
establishes what was open, cancels each order, reports the cancellation to the
member that placed it, and halts.

\begin{gap}{BUG-0048}
  Nothing bounds how far back the log is kept. The venue discards no records at
  all, so the log grows for as long as the venue runs, and what it retains is
  anchored to no position any component may present.
\end{gap}

\textbf{A checkpoint records the set itself}: the orders the component holds
open, and the sequence number those orders are current to. The region is that
checkpoint. It is the copy that outlives the process, and the log supplies only
what has happened since.

\subsubsection{How a restarted instance becomes current}
\label{sec:ha:disabled:me:recovery}

The engine reads the checkpoint and rebuilds the set as it stood at the sequence
number the checkpoint records. It then asks the sequencer for everything
committed after that number. It applies what comes back without issuing execution
reports, because those reports were issued the first time. It begins accepting
orders once it has caught up.

\textbf{A promoted instance performs the same operation.} A follower that is
promoted already holds a replica of the set and the sequence number that replica
reached. It asks for everything after that number, applies the result silently,
and begins accepting orders. The two cases differ only in where the starting
number comes from: a file after a restart, and a replica after a promotion.

Two things follow. The configuration with high availability enabled needs no
separate promotion path, because promotion is this same recovery starting from a
number learned another way. The mechanism also runs on every restart, so it is
exercised continually rather than only when a test produces a failover.


\subsubsection{What the member is told}
\label{sec:ha:disabled:me:member}

\begin{requirement}{R-0028}{Every accepted order is reported to its member, once}
  An order the matching engine accepted shall have an execution report reach the
  member that placed it, whatever restarted between the acceptance and the
  delivery, and the member shall not be led to believe the acceptance happened
  twice.
  \rationale{Catching up after a restart applies orders without issuing reports,
  on the ground that the reports were issued the first time. An order accepted
  whose report never left the engine is exactly the case where that ground does
  not hold, so something must be able to tell the two apart.}
  \uncovered{no scenario restarts the engine between acceptance and delivery}
\end{requirement}

\begin{requirement}{R-0029}{A cancel request is always answered}
  A request to cancel shall receive either a report saying the order was
  cancelled or a rejection saying why, whatever restarted while the request was
  in flight.
  \rationale{A member that does not know whether its cancel took effect must assume the order
  is live, and it cannot hedge in either direction. A member can act on a
  rejection. It cannot act on silence.}
  \uncovered{no scenario restarts the engine with a cancel in flight}
\end{requirement}

\begin{requirement}{R-0030}{With no matching engine, the venue refuses rather than acknowledges}
  Where no matching engine exists to process orders, the venue shall refuse
  orders and tell members why, rather than accept orders it cannot process.
  \rationale{An order acknowledged and never processed is worse for a member
  than one refused: it believes the order live, may have hedged against it, and
  has nothing to cancel.}
  \uncovered{scenario 42 covers refusal, and not a restart}
\end{requirement}

\begin{requirement}{R-0018}{An open order survives a restart of the matching engine}
  An order the venue held open before the restart shall be held open afterwards,
  on the same terms, and shall remain cancellable by the member that placed it.
  \rationale{The member was told the order was accepted, and has no reason to think a venue
  process has restarted. An order that stops existing without any report is a
  failure the member cannot detect, cannot cancel and cannot hedge against.}
  \covers{ha_test:50}
\end{requirement}

\begin{requirement}{R-0120}{A cancel for an order the venue no longer holds says which case it is}
  Where a member asks to cancel an order the venue does not currently hold, it
  shall be told whether the order once existed and is finished, or was never
  known to the venue at all.
  \rationale{These are different facts about the member's own position and it
  cannot establish either without being told. The protocol distinguishes them for
  that reason.

  There is a second reason particular to this venue. \emph{Unknown order} is also
  what a member is told when the venue has lost what it was holding. Answering both
  cases the same way leaves a member unable to distinguish a venue that has lost
  its orders from one that is answering correctly. Distinguishing the two cases is
  what stops that failure appearing as an ordinary reply.}
  \uncovered{no scenario cancels an order that has already been cancelled}
\end{requirement}

\begin{gap}{}
  A cancel for an order the venue does not currently hold is answered as an order
  it does not recognise, whether the order was cancelled a moment ago or never
  existed. A member cannot tell those apart, and cannot tell either of them from
  the venue having lost what it held.
\end{gap}

\begin{requirement}{R-0116}{The venue recovers what it held before deciding what to do with it}
  A restarted instance shall rebuild the set of orders it held before deciding
  whether to resume trading or to cancel them.
  \rationale{Cancelling an order requires knowing which orders to cancel and which member to
  tell about each one. A venue that cancels without having recovered cannot report
  anything, and stops mentioning the orders instead, which is what \reqref{R-0020}
  exists to prevent. Recovery is therefore required for both outcomes, and it is
  what makes cancelling possible.}
  \covers{ha_test:51}
\end{requirement}

\begin{requirement}{R-0123}{A member is told what it had, even when the venue's own record of it is unusable}
  Where the record a component keeps of the orders it held cannot be used ---
  absent, damaged, or written by a different build --- the venue shall establish
  what was open from its record of what it took, cancel each order, and report
  the cancellation to the member that placed it. Where it cannot establish that
  either, it shall halt and say that it cannot account for what it was holding.
  \rationale{\reqref{R-0102} says such a record is not to be used. That leaves
  the venue unable to name what it held, and it cannot cancel what it cannot
  name --- so without a second route it would fall silent about every order,
  which is the outcome \reqref{R-0020} exists to prevent.

  The record of what the venue took is that second route. Rebuilding from it is far
  slower than reading a checkpoint, so it is not the ordinary recovery. This case
  is rare and happens once, and the time it takes buys every member a report of
  what became of its orders.

  Reporting that the venue cannot say is the last resort, and it is still better
  than silence. A member told that the venue has lost track of its orders can act
  on that. A member told nothing cannot.}
  \covers{ha_test:52}
\end{requirement}

\begin{requirement}{R-0117}{Orders open across a long absence are cancelled, reported, and trading halts}
  Where the venue was unable to match for longer than a stated period and orders
  were open when it stopped, it shall cancel each of them, send the member that
  placed it a report saying so, and halt. Where nothing was open, it shall resume
  without halting.
  \rationale{\textbf{The period is the length of the outage, not the age of the
  order.} An order that has rested since the morning in a working market is not a problem,
  because its owner could have cancelled it at any moment and chose not to. The
  hazard is that the owner was \emph{locked out of the order}: the venue could not
  match, so the member could not cancel, and the price moved while the member could
  do nothing. The order is now priced for a market that no longer exists, and its
  owner had no opportunity to act on that.

  The hazard therefore does not depend on how long the order had been resting, nor
  on whether the order was recoverable. An order recovered correctly, whose member
  was locked out of it for ten minutes, carries the same hazard as one recovered
  from a damaged region.

  Cancelling removes the uncertainty. The reports tell each member where it stands.
  The halt records that the venue did not simply carry on. The exception for an
  empty book matters because traffic is bursty: a long absence during a quiet period
  leaves no stale order and nothing to report, and halting would produce an outage
  for no benefit.

  The venue can measure both parts of this condition, the length of the absence and
  what it cancelled, so it can enter the condition without a person present.
  Lifting the condition still requires a person, under \reqref{R-0023}.}
  \covers{ha_test:51}
\end{requirement}

\begin{requirement}{R-0118}{How long an absence may be before orders are cancelled is a stated figure}
  The period in \reqref{R-0117} shall be a figure the venue states, measured from
  the moment the venue stopped being able to match to the moment it resumes, and
  the venue shall determine that without depending on anything outside the
  component concerned.
  \rationale{The figure covers the whole absence rather than the recovery step alone.
  Noticing the death, restarting the process and rebuilding the set all count
  towards it, and the restart alone was measured at about two and a half seconds. A
  figure set against recovery time is set against the wrong quantity, and is
  reached later than intended.

  \textbf{It is measured from when matching stopped, not from the last order
  matched.} Traffic comes in bursts, so a quiet period before a failure is ordinary. Taking
  the last order matched as the start would read a ten-minute quiet period followed
  by a two-second restart as a twelve-minute absence, and would cancel every
  member's orders because the market had been quiet.

  \textbf{The component must be able to establish the figure alone.} A supervisor
  records when it killed a process and when it started the replacement, so it is
  the obvious source, but nothing in this chapter may depend on a supervisor being
  present. A component started by hand, in a debugger, hours after the last one
  stopped must reach the same answer, and it does: an hour has passed, the orders
  are an hour stale, and cancelling them is correct.

  The figure is a trading decision rather than a measurement. It states how long a
  member's resting order stays meaningful in this market, which is a judgement
  about members rather than about how long the code takes.}
  \covers{ha_test:51}
\end{requirement}

\begin{requirement}{R-0036}{A cancelled order does not become open again}
  An order the member cancelled before a restart shall not be open after it.
  \rationale{Recovery reads a record of the state and then applies what followed it. A
  cancellation that falls between the two can leave the order open again. The
  member was told the order was cancelled, so it is not watching for that order.}
  \uncovered{no scenario cancels an order and then restarts the engine}
\end{requirement}

\begin{requirement}{R-0020}{An order the venue cannot recover is cancelled to its member}
  Where the set of open orders cannot be rebuilt, the venue shall send an
  unsolicited cancel for every order it can still name, and shall say that it
  cannot name the rest.
  \rationale{A member's view of its orders matches the venue's only if the venue reports what
  it has lost. Ceasing to mention an order reports nothing.

  This requirement is the fallback beneath \reqref{R-0073}, and applies where
  recovery was attempted and could not complete.}
  \uncovered{no scenario makes recovery fail and then checks what the member is sent}
\end{requirement}

\begin{requirement}{R-0021}{A member is not told an order is unknown while the venue is still recovering}
  A request concerning an order shall not be refused as unrecognised before the
  set of open orders has been brought up to date.
  \rationale{Until the set is up to date, the venue cannot establish whether it holds the
  order. Answering that the order is unrecognised makes a statement the venue
  cannot support, and the member acts on that statement.}
  \uncovered{no scenario asserts that open orders survive a restart}
\end{requirement}

\begin{requirement}{R-0108}{The work needed to bring a component current is bounded}
  A checkpoint shall be taken often enough that the records a recovering
  component must apply after reading one do not exceed a stated bound, and that
  bound shall be a figure the venue can state.
  \rationale{Recovery takes as long as reading the checkpoint plus applying everything
  recorded since it. An unstated interval between checkpoints leaves the recovery
  time unstated too, and the interval is then set by whatever was convenient rather
  than by how long the venue may be absent.}
  \uncovered{no scenario measures the records applied after a checkpoint is read, and no bound is stated to measure against}
\end{requirement}

\begin{requirement}{R-0102}{A checkpoint that cannot be trusted is not used}
  A component shall establish that a checkpoint is whole before rebuilding from
  it, and shall not act on the part it could read. A checkpoint that is absent,
  incomplete or inconsistent shall be treated as no checkpoint at all.
  \rationale{A component reads a checkpoint once, at the moment no other record of what it
  holds is available, and it acts on the whole of what it reads. A partial
  checkpoint produces a set of orders that is wrong, and no later check detects the
  error, because every subsequent check is made against that set. Treating a
  doubtful checkpoint as absent gives the outcome of \reqref{R-0020}, which is
  worse for a member but is known.}
  \covers{ha_test:52}
\end{requirement}

\begin{requirement}{R-0022}{The venue keeps enough history to bring a recovered set of orders up to date}
  The record of committed orders shall be retained at least as far back as the
  oldest checkpoint any component may restart from.
  \rationale{Recovery reads a checkpoint and then asks for what followed it. If what followed
  has been discarded, that recovery fails in the way hardest to notice: the
  checkpoint loads, the catch-up returns nothing, and the resulting set appears
  complete.}
  \uncovered{no scenario exhausts retention, which is anchored to nothing measurable}
\end{requirement}

\subsection{Orders that outlive the trading day}
\label{sec:ha:disabled:long_lived}

Most orders last for the session that placed them, or for the trading day. A
member may also place an order to rest for weeks, and expects to find that order
open each morning without doing anything.

Everything else in this section concerns state that must survive a process or a
machine. These orders must survive \emph{the venue}: every process stopped
deliberately, the day ended, the machines rebooted, and trading started again the
next morning.

\subsubsection{Why the mechanisms already described do not serve}
\label{sec:ha:disabled:long_lived:why}

A checkpoint records the orders held at a point in the sequence. It is useful
only while the record it is anchored to still holds what followed that point, as
\cref{sec:ha:disabled:me:durable} describes. A day boundary breaks that
relationship, because a new day's record does not continue the previous day's.
Retaining every day's record, so that a checkpoint from any day remains usable,
would require the venue to keep every record it has ever written.

An order intended to outlive the day therefore needs a durable store of its own.
The venue holds that store independently of the record of what it did, and reads
it when trading begins rather than recovering it when something fails.

\textbf{No component reaches that store while it is trading.} The venue's
components hold no connection to the store, for the same reason that admission
does not depend on one: where the venue needs the store to be reachable in order
to work, every outage of the store becomes a trading outage. The venue exports
the orders and reads them when trading begins. It brings the store up to date
afterwards from its own record of what it did. The store therefore lags the
venue, and the venue can bring it level again from that record.

That has one consequence for the requirement below. \emph{The venue's own
record} must be durable before the venue tells a member its order has been
cancelled. The store is brought up to date afterwards.

\subsubsection{What the venue must do}
\label{sec:ha:disabled:long_lived:behaviour}

\begin{requirement}{R-0105}{An order that outlives the trading day is held where it survives the venue stopping}
  An order whose terms say it remains open beyond the trading day shall be
  recorded where it survives every component stopping and the day ending.
  \rationale{The venue told the member that the order would rest until the member said
  otherwise. Holding the order only where the day's own state is held makes that
  true for one day, which is not what the member was told.}
  \uncovered{no scenario stops the venue and starts it again with such an order open}
\end{requirement}

\begin{requirement}{R-0106}{Such orders are open when trading next begins}
  When trading begins, every order that was open at the end of the previous day
  and has not since expired shall be open, on the terms it was placed under, and
  cancellable by the member that placed it.
  \rationale{A member that must place its resting orders again each morning does not have
  orders that outlive the day. It has orders that the venue cancels overnight
  without reporting the cancellation.}
  \uncovered{no scenario spans a trading day boundary, which the venue does not mark}
\end{requirement}

\begin{requirement}{R-0114}{The store is known to be readable before trading begins, not during it}
  Whether the store holding these orders can be read shall be established before
  trading starts, and its being unreadable shall stop trading from starting
  rather than be discovered once it has.
  \rationale{The venue reads these orders once, at the start of a day. A store that cannot be
  read at that moment leaves the venue opening with an empty set and nothing to
  indicate that anything is missing. Every member holding such an order finds it
  gone. The check costs little, and the venue must make it before it admits
  members.}
  \uncovered{no scenario starts trading with the store unreadable}
\end{requirement}

\begin{requirement}{R-0107}{A change to such an order is committed before the member is told of it}
  Where such an order is cancelled, amended or filled, that shall be committed to
  the venue's own record before the member is told, and the store the order
  outlives the day in shall be brought up to date from that record.
  \rationale{A member told that its order is cancelled stops watching for that order. Where
  the venue writes nothing durable first, a failure between the two brings the
  order back the next time the venue reads the store. The order is then live, on a
  market that has moved, for a member with no reason to look at it. Requiring the
  venue's own record rather than the store is what allows the store to lag. The
  venue can supply the store from that record with anything the store is missing,
  and it can supply nothing for a change that was never committed anywhere.}
  \uncovered{no scenario cancels such an order and then restarts the venue}
\end{requirement}

\begin{gap}{}
  Orders that outlive the trading day are held only where every other open order
  is held, which does not survive the venue stopping. A member placing one is
  given an order that lasts until the next time the venue is restarted, and is
  told nothing to that effect.
\end{gap}

\subsection{The matching engine publisher}
\label{sec:ha:disabled:mep}

The publisher takes what the venue has done and makes it available to consumers
inside the venue: market data, and the recorder described in
\cref{sec:ha:disabled:oar}. The publisher dies and is restarted. No member
notices this directly. What stops is the flow of records to those consumers.

\subsubsection{What state it holds}
\label{sec:ha:disabled:mep:state}

The publisher holds the record of what it has been given, and how far each
subscriber has consumed that record. It does not have to retain the second of
those across a restart. A subscriber presents its own position when it
reconnects, so the publisher is told where each subscriber has reached.

\subsubsection{Where the copy that outlives it lives}
\label{sec:ha:disabled:mep:durable}

The copy lives in the publisher's own log, on disk. The slowest subscriber that
has not yet consumed a record determines what the publisher may discard from that
log.

\begin{gap}{BUG-0048}
  What the publisher may discard is anchored to no position at all. It discards
  nothing, so its log grows for as long as the venue runs, and the positions its
  subscribers have reached are never what decides that a record is kept.
\end{gap}

\subsubsection{How a restarted instance becomes current}
\label{sec:ha:disabled:mep:recovery}

The publisher reopens its log and waits for each returning subscriber to present
its position. The publisher has one such operation rather than two, as the
gateway in \cref{sec:ha:disabled:gateway} does.

\subsubsection{What the member is told}
\label{sec:ha:disabled:mep:member}

Nothing. A member is not a subscriber. The requirements in this subsection
protect the consumers instead.

\begin{requirement}{R-0047}{A subscriber resumes where it had reached, across a restart of the publisher}
  The venue shall supply a returning subscriber with every record after the
  position it presents, and shall not require it to have been present
  throughout.
  \rationale{A subscriber that must stay connected continuously to avoid missing records
  cannot be restarted, upgraded, or briefly separated from the network. No operator
  could run such a subscriber.}
  \uncovered{no scenario restarts the publisher under a connected subscriber}
\end{requirement}

\begin{requirement}{R-0048}{A record is retained until every subscriber has consumed it}
  The venue shall not discard a record that any subscriber entitled to it has
  not yet consumed.
  \rationale{Where the publisher discards records on a timer rather than when they have been
  consumed, a subscriber that was briefly away loses records and nothing reports
  the loss. The subscriber discovers the loss, and the publisher that caused it
  does not.}
  \uncovered{no scenario returns a subscriber after its records were discarded}
\end{requirement}

\subsection{The order activity recorder}
\label{sec:ha:disabled:oar}

The recorder subscribes to the publisher and forwards what the venue has done to
an enterprise bus, which is how systems outside the venue learn of it. A failure
of the recorder is felt outside the venue rather than within it.

\subsubsection{What state it holds}
\label{sec:ha:disabled:oar:state}

The recorder holds one number: how far it has published and had confirmed. It can
obtain everything else again from the publisher.

\subsubsection{Where the copy that outlives it lives}
\label{sec:ha:disabled:oar:durable}

The copy is the recorder's own record of that number, written to disk. The
recorder \textbf{writes it only after the bus has confirmed the publish, and
never before}. That ordering is what makes the design durable. A failure between
publishing and recording the position causes the recorder to publish the record
again, and the consumers downstream handle the duplicate. A failure in the other
order skips a record permanently, and nothing afterwards can detect the omission.

The recorded position therefore needs to be no more than a lower bound. The
recorder can write it as late as it likes, whatever the position's staleness,
because a stale position costs duplicate records and nothing else.

\subsubsection{How a restarted instance becomes current}
\label{sec:ha:disabled:oar:recovery}

The recorder reads its position, presents it to the publisher, and continues. It
publishes again any record it had already published but not recorded. This
duplication is expected: a consumer of the enterprise bus is required to cope
with records it has already seen.

\subsubsection{What the member is told}
\label{sec:ha:disabled:oar:member}

Nothing directly. The last requirement in this subsection is the exception,
because it halts trading.

\begin{requirement}{R-0049}{Nothing is lost from the external record, though it may appear more than once}
  Every record the venue produces shall reach the enterprise bus at least once.
  Duplicates are permitted; omissions are not.
  \rationale{The venue cannot repair a loss outside the venue, because nothing outside it
  knows which records are missing. Anybody holding the record's identity can repair
  a duplicate.}
  \uncovered{no scenario, and none is possible until the recorder exists}
\end{requirement}

\begin{requirement}{R-0050}{Every externally published record carries an identity a consumer can deduplicate on}
  Each record shall carry a stable producer identity and a sequence number
  assigned by the venue.
  \rationale{This requirement is what makes \reqref{R-0049} usable. Without an identity, a
  consumer cannot distinguish a duplicate from a second event, and permitting
  duplicates would then permit the record to be wrong.}
  \uncovered{no scenario, and none is possible until the recorder exists}
\end{requirement}

\begin{requirement}{R-0051}{The recorder never resumes beyond its last confirmed publish}
  The position the recorder resumes from shall never be later than the last
  publish the bus confirmed.
  \rationale{Recording a position before the bus confirms the publish turns a crash into a
  permanent omission, which \reqref{R-0049} forbids. Relaxing the ordering to
  increase throughput would reintroduce that omission.}
  \uncovered{no scenario, and none is possible until the recorder exists}
\end{requirement}

\begin{requirement}{R-0052}{A restart of the recorder does not stop trading}
  The venue shall continue to accept orders while the recorder restarts and
  catches up.
  \rationale{A supervised restart is a routine event, and the venue recovers from it without
  help. Halting trading for such a restart would make the recorder's availability,
  rather than the venue's, decide whether members can trade.}
  \uncovered{no scenario, and none is possible until the recorder exists}
\end{requirement}

\begin{requirement}{R-0053}{Where external publishing cannot resume, trading is halted}
  Where no recorder is publishing and the unpublished backlog has passed a
  stated bound, the venue shall halt trading, and shall say that it has.
  \rationale{Nobody can reconstruct trading afterwards once it has outrun the record of it.
  The venue can measure both parts of this condition: whether a recorder is
  connected, and how large the backlog is. It can therefore act on the condition
  without a person present. The bound must exceed the time a supervised restart and
  catch-up take, or it would contradict \reqref{R-0052}, so the bound has to be
  measured rather than chosen.}
  \uncovered{no scenario, and none is possible until the recorder exists}
\end{requirement}

\subsection{The sequencer and the matching engine together}
\label{sec:ha:disabled:seq_and_me}

Both processes die, and both are restarted. This case is not the two preceding
subsections combined, because each of those recoveries relies on a component that
has also died here.

\subsubsection{What state it holds}
\label{sec:ha:disabled:seq_and_me:state}

Between them, the two components hold everything the venue has about work in
progress. The sequencer holds the record of what has been taken. The matching
engine holds which orders are open. Neither holds the other's state, and neither
can reconstruct it.

\subsubsection{Where the copy that outlives it lives}
\label{sec:ha:disabled:seq_and_me:durable}

Two records survive. The sequencer's log holds what the venue took and the order
it took them in. The matching engine's region holds the orders it had open and
the position those orders are current to, which
\cref{sec:ha:disabled:me:durable} describes.

What does not survive is the sequencer's record that particular orders were
waiting for an engine. Only the running process holds that.

This case differs from the single failures because the component each of those
recoveries relies on has also died. A matching engine becomes current by asking
the sequencer, so it cannot become current until the sequencer is current
itself.

\subsubsection{How a restarted instance becomes current}
\label{sec:ha:disabled:seq_and_me:recovery}

This subsection is about an order that sits between the two components:
committed by the sequencer, and never given to the matching engine.

The engine establishes the position its region is current to, asks the sequencer
for everything committed after that position, and applies what comes back. The
order is in the log, so the sequencer serves it. The engine applies it and sends
the member an execution report for it, marked as a report the member may already
have seen, which is \reqref{R-0122}.

The sequencer's own record that the order was waiting does not survive the
restart, and the venue does not depend on it. The position the engine presents is
what determines the records the sequencer sends.

Recovery therefore acquires an ordering it does not otherwise have. The engine
cannot become current before the sequencer is current, because the engine becomes
current by asking the sequencer. Where only one of the two components failed, the
surviving component was already running and that ordering did not arise.

\subsubsection{What the member is told}
\label{sec:ha:disabled:seq_and_me:member}

\begin{requirement}{R-0015}{Trading does not resume while committed orders remain unresolved}
  The venue shall not accept new orders after such a failure until every order
  it had committed has either been given to the matching engine or had an
  execution report produced for it. Whether a member is present to receive that
  report does not bear on this: a member that is slow to return, or that never
  returns, cannot hold the venue shut.
  \rationale{Accepting new orders while earlier orders remain unresolved conceals the loss
  rather than recovering from it. The venue then looks healthy to a member whose
  earlier order the venue will never mention again.}
  \uncovered{no scenario restarts both components together}
\end{requirement}

Where the automatic paths cannot resolve a failure, a person resolves it, and a
longer timer does not. The venue therefore needs a halted state that a person can
put it into. The venue also needs that state for reasons unconnected to any
failure: a scheduled pause, a market-wide event, or an instrument trading in
disorderly conditions.

\begin{requirement}{R-0023}{Trading can be halted by a person, and members are told}
  An operator shall be able to halt trading. Members shall be told that the
  venue has halted and shall be told when it resumes, and the halt shall not
  lift of its own accord.
  \rationale{A member that is connected but idle sends no order, so the rejection of an order
  tells it nothing. A member reconnecting into a halt would otherwise discover the
  halt by sending an order and having it refused. Requiring a person to lift the
  halt is what distinguishes it from the automatic states that clear themselves.}
  \uncovered{no scenario halts the venue, because it has no halted state to enter}
\end{requirement}

\begin{requirement}{R-0024}{A halt that cancels open orders reports every cancellation}
  Where a halt cancels the venue's open orders, an execution report shall be
  sent to the member that placed each one. No order shall cease to be open
  without the member that placed it being told.
  \rationale{Cancelling the open orders is a normal thing for a halt to do, and real venues
  have done it when a market became disorderly. What makes it safe is that each
  member's view of its orders comes to match the venue's, and that requires the
  venue to report what it cancelled rather than to stop mentioning those orders.}
  \uncovered{no scenario halts the venue, because it has no halted state to enter}
\end{requirement}

\begin{requirement}{R-0017}{Recovery does not depend on the order the components return in}
  The venue shall reach the same state whichever of the sequencer and the
  matching engine starts first after both have failed.
  \rationale{A launcher restarts what it starts and coordinates with nothing, so the venue
  does not decide which of the two components returns first. A component that
  cannot yet be told what it needs must wait, rather than proceed on the assumption
  that there is nothing to receive.}
  \uncovered{no scenario restarts both components together}
\end{requirement}

\subsection{The gateway together with another component}
\label{sec:ha:disabled:gw_and_others}

The gateway depends on the sequencer for everything it recovers, and on the
matching engine for everything it asks to be done on a member's behalf. Losing
either of those alongside the gateway is therefore a separate case. This
subsection states each case even where the requirements are stated elsewhere, so
that a reader can see which combinations the venue is designed to survive.

\subsubsection{The gateway and the sequencer}
\label{sec:ha:disabled:gw_and_seq}

A member reconnects to a restarted gateway, and that gateway asks a sequencer
that has itself just restarted and no longer holds the session. This case needs
no new requirement. \reqref{R-0011} requires the numbering to survive, and
\reqref{R-0012} requires the venue to say so where it has not, rather than
resuming the member at the beginning. The combination is the likeliest route to
\reqref{R-0012}, because no component reports a session that was not bound at the
moment the sequencer restarted.

\subsubsection{The gateway and the matching engine}
\label{sec:ha:disabled:gw_and_me}

A member's connection ends while the matching engine is not running. The gateway
does not act on the member's orders itself. It reports the disconnection and the
terms it was given to the component that holds the orders, and that component is
absent.

\begin{requirement}{R-0037}{A member's cancel-on-disconnect instruction is not lost because a component was restarting}
  Where a session ends while the component that holds its orders is unavailable,
  the instruction shall be acted on once that component returns, or the member
  shall be told that it was not.
  \rationale{The member configured this to bound its exposure when its own connection fails.
  An instruction dropped because a venue process was restarting leaves the member
  exposed, and neither the member nor the venue holds any record of it.}
  \uncovered{no scenario ends a session while the engine is down}
\end{requirement}

\subsection{The member's own application}
\label{sec:ha:disabled:client}

A member's application dies and is restarted. This is the only failure in this
section that happens outside the venue, so the venue's obligations differ. The
venue has nothing to recover here, and must hold everything it was keeping for
that member.

It is also the failure the venue meets most often. A member restarts its systems
on its own schedule, for its own reasons, and tells the venue nothing beforehand.

\subsubsection{What state it holds}
\label{sec:ha:disabled:client:state}

The member's application holds its own view: which orders it treats as live at
the venue, the identifiers it has used, and how far its side of the session
numbering had reached. The venue does not keep any of this, and must not assume
that any of it survived.

The venue holds a different set of state \emph{on the member's behalf}, and the
rest of this subsection concerns it: the session's numbering, the orders open
under that session, and the terms the member was admitted under.

\subsubsection{Where the copy that outlives it lives}
\label{sec:ha:disabled:client:durable}

The member chooses where to keep its own copy. The venue holds no record of that
choice and must not depend on it. A member may return holding its full position,
or holding none of it, and both are legitimate.

The venue must therefore answer the member rather than assume anything about it.
A returning member that has kept its numbering resumes at that numbering. A
member that has lost its numbering says so and starts again from the beginning.
In both cases the member must be able to establish what the venue holds for it,
which \cref{sec:ha:disabled:gateway:member} records that it currently cannot.

\subsubsection{How a restarted instance becomes current}
\label{sec:ha:disabled:client:recovery}

The member's application connects, logs on, and either resumes its numbering or
asks to restart it. The venue supplies what it holds for the session, and the
member reconciles its own view against what the venue reports.

\subsubsection{What the member is told}
\label{sec:ha:disabled:client:member}

\begin{requirement}{R-0054}{A member that returns in time finds its orders untouched}
  Where a member's connection ends unexpectedly and the same comp id logs on
  again within the period provisioned for it, none of that member's orders shall
  have been cancelled.
  \rationale{The period exists so that a member's own restart is a reconnection rather than an
  event the market sees. Cancelling the orders and having the member place them
  again costs the member its position in the queue, and reports activity to the
  market that did not happen.}
  \covers{ha_test:19}
\end{requirement}

\begin{requirement}{R-0055}{A member that does not return has its orders cancelled as it asked, and is told when it comes back}
  Where the period expires without the member returning, the orders it asked to
  have cancelled shall be cancelled, and the reports of those cancellations
  shall be delivered to it when it next connects.
  \rationale{The member configured this to bound its exposure while it is absent. Cancelling
  the orders without keeping the reports for the member's return leaves the member
  to work out what happened for itself.}
  \uncovered{no scenario lets the period expire and then reconnects}
\end{requirement}

\begin{requirement}{R-0056}{An order meant to outlive the session is not cancelled by a disconnection}
  An order whose time in force says it survives the trading session shall not be
  cancelled because the member's connection ended.
  \rationale{The member placed such an order to outlive the session that placed it.
  Cancelling it because a connection closed does the opposite of what the member
  asked for, and the member cannot restate the instruction cheaply.}
  \uncovered{no scenario places an order that outlives the session and then drops it}
\end{requirement}

\begin{requirement}{R-0057}{A member that has lost its own numbering can start afresh without losing its orders}
  A member that asks to restart its numbering shall be given a session numbered
  from the beginning, and the orders it holds at the venue shall be unaffected.
  \rationale{Losing a sequence store is an ordinary consequence of a member's own restart, and
  it says nothing about that member's orders. Tying the two together would make a
  member choose between a session it can use and the orders it holds.}
  \uncovered{no scenario restarts a member's numbering with orders open}
\end{requirement}

\subsection{Supervision}
\label{sec:ha:disabled:supervision}

A launcher answers one question: is this component running, and if not, start it.
It takes no part in deciding roles, and nothing in this chapter depends on it
being present. This subsection states the bounds on what it must do.

\begin{requirement}{R-0033}{A component asked to stop is not restarted}
  Where a component is stopped deliberately, whatever supervises it shall not
  start it again.
  \rationale{Without this requirement, nobody can stop the venue. Every stop becomes a
  restart, and an operator has no way to take a component out of service.}
  \uncovered{no scenario stops a supervised component deliberately}
\end{requirement}

\begin{requirement}{R-0038}{A component that cannot start at all is reported at once}
  Where a component cannot be started --- its binary absent, its configuration
  unreadable, a port it needs already held --- whatever supervises it shall
  report the failure immediately and shall not retry as though the attempt might
  succeed.
  \rationale{Retrying a start that cannot succeed turns a fault an operator could fix in a
  minute into an absence that nothing reports. None of these causes is transient,
  so waiting achieves nothing. A component that starts and then dies is a different
  case, and \reqref{R-0039} covers it.}
  \uncovered{no scenario starts a component that cannot start}
\end{requirement}

\begin{requirement}{R-0039}{A component that fails repeatedly and quickly is not restarted again}
  Where a component dies repeatedly within a short time of being started,
  whatever supervises it shall stop restarting it. A death soon after a start is
  a failure to start rather than a run that ended, and shall be counted as one.
  \rationale{Restarting a component whose fault is deterministic produces another failure, and
  restarting it without limit produces a fault that runs continuously and reports
  nothing. Distinguishing a failure to start from a run that ended is what makes
  the condition detectable: a healthy component runs for hours, and one that exits
  in seconds has not run.}
  \uncovered{no scenario makes a component fail repeatedly}
\end{requirement}

\begin{requirement}{R-0040}{Giving up is announced where something other than a person can see it}
  Having stopped restarting a component, the supervisor shall say so on its
  output, in a form a log monitor can detect and raise an alert from without
  anyone watching at the time.
  \rationale{Automatic restart removes the symptom that would otherwise draw attention.
  Somebody downstream eventually notices a component that dies and stays dead.
  Nobody notices a component that is restarted forty times an hour. Automating the
  recovery therefore makes the reporting necessary.}
  \uncovered{nothing asserts what the supervisor says when it gives up}
\end{requirement}

\begin{requirement}{R-0041}{A component that runs healthily again is given a fresh allowance}
  Where a restarted component runs for long enough to have started
  successfully, the count of rapid failures against it shall be reset.
  \rationale{Without this requirement, the allowance is consumed over the life of the
  deployment rather than over one episode. A component that fails once a month
  eventually exhausts the allowance, and the supervisor then stops restarting it
  for a fault it has recovered from every time.}
  \uncovered{no scenario makes a component fail repeatedly}
\end{requirement}

\subsection{Stopping and starting deliberately}
\label{sec:ha:disabled:stopping}

Not every stop is a failure. An operator takes a component down for a release,
for a configuration change, or because the trading day has ended. None of those
should cost the venue what a crash costs.

\begin{requirement}{R-0034}{An orderly stop leaves less to recover than a crash}
  A component asked to stop shall record its state before it goes, so that what
  it has to recover afterwards is no more than what it wrote.
  \rationale{A planned restart is the common case, for a release or a configuration change,
  and it need not cost what an unplanned restart does. It is also the cheapest way
  to exercise the recovery path regularly.}
  \uncovered{no scenario stops a component deliberately, and nothing would be written if one did}
\end{requirement}

\begin{requirement}{R-0035}{The start and end of a trading day are published}
  The venue shall publish the start of a trading day and its end as events that
  any interested party may subscribe to.
  \rationale{A venue that has just started is either resuming a trading day or beginning a new
  one, and which of those it is changes what every recovery in this section means.
  A component must be able to ask rather than infer it. Publishing the event once
  to every subscriber also stops each component being told separately, which is how
  components come to hold different answers.}
  \uncovered{no scenario spans a trading day boundary, which the venue does not mark}
\end{requirement}


\section{Recovery is one mechanism}
\label{sec:ha:recovery_shape}

\Cref{sec:ha:disabled} asks the same four questions of every component. The
answers to the third question, how a restarted instance becomes current, share a
common shape. This section sets that shape out, because the next section depends
on it.

\textbf{Every component recovers the same way.} Whatever a component holds, and
whichever component holds the authoritative record of it, the operation has four
steps:

\begin{enumerate}[nosep]
  \item establish a \emph{position}: a point up to which its state is current;
  \item present that position to whatever holds the authoritative record;
  \item receive everything after it and apply that silently, without repeating
        side effects that were performed the first time;
  \item begin acting.
\end{enumerate}

What differs between components is only where the position in step one comes
from.

\begin{center}
\begin{tabular}{lll}
\toprule
\textbf{Component} & \textbf{Position, on a restart} & \textbf{Position, on a promotion} \\
\midrule
Authentication service & how current its credentials are & no promotion \\
Gateway                & the sequencer, on binding       & no promotion \\
Sequencer              & its log, and gateways reporting & the replica it maintained \\
Matching engine        & a checkpoint on disk            & the replica it maintained \\
Publisher              & each subscriber presents its own & each subscriber presents its own \\
Recorder               & its own confirmed publish        & its own confirmed publish \\
\bottomrule
\end{tabular}
\end{center}

Three things follow from this.

\textbf{Promotion is not a second mechanism.} For the components that have one,
promotion is this same operation with the position learned from a replica rather
than from a file or from being told. \Cref{sec:ha:enabled} is therefore not
another recovery design; it is about how an instance learns it may act, and about
the failures that exist only because there are two of them.

\textbf{The venue exercises this mechanism constantly rather than rarely.} A
recovery path used only on promotion runs only when a scenario produces a
promotion. The same path, used on every restart, runs whenever a component is
upgraded, reconfigured or dies. The venue therefore has evidence that the path
works.

\textbf{The requirements are shared, and this document states them once.} The
operation is common to every component, so the properties it needs are common
too, and no component restates them: \reqref{R-0031}, that no component acts on
partially recovered state; \reqref{R-0032}, that an interrupted recovery can be
repeated without applying anything twice; and \reqref{R-0064}, that the
generation stamped on what is served never goes backwards.

\textbf{Every authoritative source carries the same retention obligation.} A
position is useful only while the records after it still exist. Each source, the
sequencer's log, the publisher's log and a member's reports, must therefore
retain records at least as far back as the oldest position any party may still
present to it. \reqref{R-0008}, \reqref{R-0022} and \reqref{R-0048} state that
bound for each source in turn. In each case the retained period must be a figure
the venue states, rather than whatever the available storage happens to hold.

\begin{requirement}{R-0101}{A component establishes that its catch-up was complete before it acts}
  A component shall establish that every record the authority sent it arrived,
  exactly once and in the order sent, and shall not begin acting if it did not.
  \rationale{A component applies catch-up records without producing the execution reports it
  produced the first time, so applying them produces no visible output. A catch-up
  that is missing a record therefore leaves a component whose state is wrong and
  that treats itself as current, and every check it performs afterwards is made
  against that wrong state. Applying a record silently is correct. Failing to
  receive one silently is not.

  \textbf{The wording states what the authority sent, rather than what lies between
  the two positions, because a component cannot establish the second.} The stream
  is filtered. An authority bringing a component back up to date sends the records
  that component consumes and withholds the records it produces, so the records
  that arrive are a subset of the numbers in the range, and the gaps between them
  are ordinary. A check for a run of consecutive numbers is therefore wrong,
  although the numbering makes such a check look natural.

  The check that remains is a count, an ordering, and an absence of repeats. It
  catches what goes wrong between two parties: loss in transit, reordering and
  duplication. \textbf{It does not catch the authority sending the wrong set},
  because a component cannot establish that it received a record it was never told
  about. This paragraph states that limit, because the requirement otherwise reads
  as a stronger guarantee than the venue provides.}
  \uncovered{no scenario withholds a record during a catch-up, so the check is
  exercised only by its unit tests and never against a real stream}
\end{requirement}

\textbf{Three components depart from this shape.} A gateway and an authentication
service are never promoted, so each has one operation instead of two. Neither
therefore needs an arbiter, and running several of each costs little. A restarted
sequencer recovers \emph{less} than a promoted one, because no component reports
a session that was not bound at the time of the restart. That difference is the
reason for \reqref{R-0012}.

\section{The venue says what it is doing}
\label{sec:ha:technical_event}

\emph{Placeholder, created 2026-09-08. The requirements below are named, not written.}

The venue can say one thing about itself --- whether it is accepting orders --- and everything
else it knows is private to the process that knows it. A matching engine that cannot account for
what it was holding sets a local flag and writes a line to its log, and no gateway and no
sequencer mentions a halt at all: the gateway goes on admitting orders and the sequencer goes on
sequencing them into a venue whose engine has stopped. ``Trading is halted'' is one
process's private opinion.

That is why this section is in this chapter rather than beside the components that will carry the
message. \Cref{sec:ha:enabled:process_death} requires that the venue not run with no instance of a
component without saying so, and it cannot be built until the venue is able to announce anything
about itself at all. The announcement is the missing mechanism, and the phases are what it
announces.

The shape is settled and is not to be re-argued here; it is recorded in
\texttt{docs/venue/trading\_phases.md}. This section turns it into requirements. In outline: the
venue is always in exactly one of five phases, of which start of day and end of day are two,
because work happens during both and a period with duration is a phase and not a moment; a halt is
a flag alongside the phase rather than a phase of its own, so that lifting one does not depend on
remembered state; and the message carries the resulting state rather than the transition, together
with the trading day, so that a component which missed an event is not left holding a private
state machine that is wrong in a way it cannot detect.

\subsection{The requirements to be written}
\label{sec:ha:technical_event:todo}

\begin{description}
  \item[The venue declares its phase, and no component derives it.] A component holds what it was
    told and never computes the phase, or the trading day, from its own clock.
  \item[An event carrying a trading day older than the one a component holds is refused.] The
    trading day plays here the part the epoch plays in leadership, so that a halt declared
    yesterday cannot stop trading today.
  \item[A component that has just started can ask rather than wait.] The same property the
    order-acceptance broadcast already provides by repeating itself while the venue is degraded.
  \item[Both audiences are served, and the two paths are named.] A component that must act on a
    phase takes it on the control plane; one that only needs to know takes it by publish and
    subscribe. A phase that some components acting on it have and others do not is worse than no
    phase at all.
\end{description}

\textbf{Two things are deliberately not settled yet}, and no requirement here is to assume either.
What each component may do in each phase is a table that does not exist --- whether a member may
cancel a resting order after the close is a policy decision with consequences for every order the
venue holds overnight, and naming the phases does not settle it. And whether a phase belongs to
the venue or to an instrument is open: this venue is single-instrument in effect, so venue-wide is
right today, but the shape of the message decides whether per-instrument is an addition later or a
change.

\section{High availability enabled}
\label{sec:ha:enabled}

This is the configuration a venue runs in. The switch starts a second instance
of every component that must be singular --- the sequencer, the matching engine
and the recorder --- together with the components whose only purpose is to
decide which of a pair may act: a pair of arbiters and a witness.

The switch changes nothing about the replicated peers. The authentication
services and the gateways run the same way in both configurations, for the reason
\cref{sec:ha:disabled:scope} gives: several of each run at once, none of them
leads, and an arbiter has nothing to decide about them.

\subsection{Two pairs of words, and why they must not be confused}
\label{sec:ha:enabled:vocabulary}

\textbf{Primary and secondary are permanent names.} They are fixed in
configuration when an instance is deployed and never change for the life of that
deployment. They say which instance is preferred when both start together and
neither has any other claim, and they say nothing else. They are not a statement
about which instance is doing the work.

\textbf{Leader and follower are positions.} The leader is the instance currently
entitled to act; the follower is the one that is not. Either configured instance
may hold either position, and which it holds may change repeatedly over a
deployment's life.

\textbf{``The primary is the follower'' is therefore not a contradiction.} After
one failover it is the ordinary state, and it remains so until something moves
leadership again. A design that treats the two pairs of words as synonyms moves
leadership back to the primary for no reason beyond its name, and moves it to the
instance that has just failed.

The pair must also be able to exchange positions repeatedly. If the secondary is
promoted and later dies, the primary, which is the follower by then, must be able
to take over again. A scheme that can fail over only once postpones an outage
instead of preventing it.

This document uses neither \emph{active} nor \emph{standby}. Readers take those
words to mean the configured identity in some places and the position in others,
which is the ambiguity this section exists to prevent.

\subsection{What changes, and what does not}
\label{sec:ha:enabled:what_changes}

\textbf{What changes} is that the venue survives the loss of a machine. The
configuration in \cref{sec:ha:disabled} can only restart what died and rebuild
what it held. This configuration has a second instance that already holds it.

\textbf{What does not change is the recovery itself.} For the matching engine,
becoming current after a restart and becoming current after a promotion are the
same operation. They differ only in where the starting point comes from: a file
after a restart, and a replica after a promotion. This section therefore
describes how an instance learns that it may act, and the failures that exist
only because two instances run.

\textbf{What is new} follows from a second instance existing: deciding which of
the two may act, ensuring that only one of them acts, detecting that the other
has gone, and refusing what an instance sends after its entitlement has passed to
the other. The subsections that follow cover those, and the components that exist
to arbitrate.

\subsection{What running two of everything has to be worth}
\label{sec:ha:enabled:worth}

Running two instances of a component adds ways for the venue to fail that running
one instance does not have. A single instance cannot disagree with itself, cannot
be separated from itself by a network failure, and cannot both act and fail to
act. A pair can reach all three conditions.

Two things must therefore hold for the second instance to be an improvement.

\textbf{The venue must not be stopped by something a single instance would have
survived.} If a pair of sequencers stops trading because they cannot reach each
other, when one sequencer alone would have carried on, the redundancy has made
the venue less available rather than more.

\textbf{The venue must prevent the new failures rather than make them unlikely.}
Two instances acting at once, and neither instance acting at all, are conditions
that cannot arise with one instance. They exist only because a second instance
exists. Making them rare is not sufficient, because the venue goes on looking
healthy in both. Nobody sees the first until the two instances' states are
compared, and nobody sees the second until somebody asks why nothing is trading.

Most of this section is about preventing them.

\subsection{The arbiter and the witness}
\label{sec:ha:enabled:arbiter}

Two instances of a group may act, and only one of them may act at a time. This
subsection says how the venue decides which one.

An instance cannot make this decision alone. When an instance loses contact with
its peer, it cannot distinguish a peer that has stopped from a peer it can no
longer reach. Both readings lead to a wrong action. If it promotes itself and the
peer is alive, both instances act and their state diverges. If it declines and
the peer is dead, neither instance acts and the venue stops trading. The second
failure is the safer one, and it takes longer to notice, because no process
crashes and nothing logs an error.

\textbf{A majority of three voters decides.} Each group has three voters: its two
instances and the arbiter pool. An instance acts only while a majority of the
three has granted it a lease that has not run out. Its own vote is one of the
three, so it needs a grant from one other voter: its peer, or the arbiter pool. A
voter grants a lease to one instance at a time, and promises not to grant another
until the lease period has passed. Two instances would each need two of the three
votes at the same moment, so at least one voter would have to have granted both,
and no voter does. Every message also carries the generation number of the
entitlement it was sent under, and a receiver refuses a message whose generation
has been superseded.

\textbf{Two arbiters run, and one witness runs alongside them.} The arbiter pool
votes through whichever arbiter is active. The two arbiters decide which of them
is active by the same rule, with the witness as their third voter: an arbiter is
active only while a majority of the two arbiters and the witness has granted it a
lease. Losing one arbiter moves no entitlement and interrupts nothing, which is
\reqref{R-0063}. The witness has no pair and no standby.
\Cref{sec:ha:enabled:arbiter:witness} states what the venue does about that.

\subsubsection{What is done about the witness being single}
\label{sec:ha:enabled:arbiter:witness}

\textbf{The witness holds no state, so losing it loses nothing.} It has no record
to recover and no history to lose. After it restarts it grants nothing for one
lease period, which is the whole of its recovery.

\textbf{Losing the witness alone changes nothing.} Two arbiters that can see each
other grant each other leases without it. The absence of the witness is felt only
when an arbiter is also lost: the surviving arbiter then holds only its own vote,
and it is active only while its existing lease lasts. The pool then has no active
arbiter until another voter returns. Component leaders are unaffected while their
peers are running, because each renews its lease with its peer.

\textbf{The witness is off every ordinary path.} Only the arbiters talk to it. No
order, report or component touches the witness.

\textbf{The witness runs where the arbiters do not.} A witness that shares a
failure domain with one arbiter gives little protection, because a single event
then removes that arbiter and the witness together, and the surviving arbiter
holds only its own vote. The witness therefore runs on a third machine, on a
different power supply, behind a different network switch, and ideally in a
different rack. Its usefulness comes from where it is deployed rather than from
what it does.

Three properties of the witness process make that placement cheap:

\begin{itemize}[nosep]
  \item \textbf{It needs no storage.} It holds nothing on disk, so it depends on
        no array and no filesystem.
  \item \textbf{It is off the hot path.} It needs no pinned cores, no low-latency
        hardware and no proximity to any other component, so it can run on the
        least capable machine in the deployment.
  \item \textbf{Its network route matters more than its speed.} The failure that
        hurts the venue is an arbiter being unable to reach the witness. The
        route to the witness should therefore fail independently of the route
        between the two arbiters.
\end{itemize}

The venue goes on trading when both arbiters and the witness are gone, for as long
as each group's leader can renew its lease with its peer. What it cannot do then is
replace a leader that fails: \cref{sec:ha:enabled:degraded} says what happens.

\subsubsection{What state it holds}
\label{sec:ha:enabled:arbiter:state}

The arbiter holds whether it is active and, while it is, one vote per group: the
instance it has promised its vote to, if any, and until when. It also holds the
highest generation granted in each group. It holds nothing about orders, sessions
or sequence numbers. That limit is deliberate: a component that decides which
instance may act must not itself hold state that can be wrong.

The witness holds its vote on which arbiter is active, and nothing else.

\subsubsection{Where the copy that outlives it lives}
\label{sec:ha:enabled:arbiter:durable}

The arbiter keeps nothing across its own restart. It does not need to know what it
promised before it stopped: it grants nothing for one lease period after starting,
by which time every such promise has run out. The same applies to an arbiter that
becomes active in place of the other, because it does not know what the other
promised.

The highest generation granted in each group is the one thing worth keeping, and
the active arbiter tells the passive one each time it rises. An arbiter that
becomes active therefore knows it, unless both arbiters have restarted.

\subsubsection{How a restarted instance becomes current}
\label{sec:ha:enabled:arbiter:recovery}

A restarted arbiter waits one lease period and then votes. Until then it refuses
every request, which is \reqref{R-0059}.

\subsubsection{What the member is told}
\label{sec:ha:enabled:arbiter:member}

Nothing. The arbiter sits off the path a member's orders travel, and a member
cannot tell that one exists. An arbiter's failures reach a member only through
the components that consult it, so the requirements here are stated in terms of
those components rather than of members.

\begin{requirement}{R-0058}{Two instances of a group are never both entitled to act}
  At any moment, at most one instance of a group shall hold the entitlement to
  act.
  \rationale{This is the property the arbiter exists to provide, and the rest of
  this subsection follows from it. Two instances acting on the same state
  diverge, and nothing afterwards can establish which of them was right.}
  \uncovered{no scenario asserts the property directly, only particular routes to it}
\end{requirement}

\begin{requirement}{R-0146}{An instance acts only while a majority of its three voters grants it a lease}
  An instance of a group shall act only while it holds a lease, not yet run out,
  from its peer or from the arbiter pool. It shall count each lease from the
  moment it asked for it, less an allowance for clocks that run at slightly
  different rates. A voter shall grant a lease to one instance at a time, shall
  grant nothing for one lease period after it starts unless it has kept a record,
  made during the same boot of its machine, of what it promised, and shall grant
  no generation below one it has already granted. An instance that has granted its
  peer a lease shall not ask to lead until that lease has run out.
  \rationale{An instance that loses contact with its peer cannot tell a dead peer
  from an unreachable one, so no rule it applies alone can be right in both
  cases. A majority of three can be: two instances would each need two of the
  three votes at the same moment, which would need a voter that has granted both.

  Each clause closes a way in which two instances come to act at once, and model
  checking showed each one is needed: without it, two instances act as leader at
  once within a few steps. Counting the lease from the request rather than the
  grant means the holder always believes its lease ends first. Waiting after a
  start covers a voter that has forgotten what it promised; a voter that kept a
  record has not forgotten, and so can resume at once, which is what lets a
  supervised restart keep the lead. The epoch rule stops
  a voter granting an older generation. An instance that granted its peer a lease
  must not then count its own vote for itself.}
  \covers{ha_test:8, ha_test:15, ha_test:39, ha_test:49, tla:lease-1-holder-counts-from-arrival, tla:lease-2-candidate-ignores-own-grant,
  tla:lease-3-restart-does-not-wait}
\end{requirement}

\begin{requirement}{R-0059}{An arbiter that does not know what was promised grants nothing until every promise has run out}
  An arbiter that has just started, or has just become the active arbiter, shall
  grant no lease to any component for one lease period.
  \rationale{An arbiter keeps nothing across its own restart, and an arbiter that
  becomes active does not know what the previously active one promised. During
  that period it cannot distinguish two situations: a group in which no instance
  holds a lease, and a group whose leader holds one it cannot see. Granting a
  lease is correct in the first situation. In the second it gives a group two
  instances entitled to act, which is the outcome the arbiter exists to prevent.

  Waiting one lease period settles it without any knowledge: by then every
  promise made before is certain to have run out. The leader of each group
  renews with its peer meanwhile, so the wait costs nothing while a pair has both
  instances running.}
  \uncovered{no scenario asks a freshly started arbiter about a group it has not been told of}
\end{requirement}

\begin{requirement}{R-0060}{Silence from an arbiter is not read as the absence of one}
  A component that asks and hears nothing shall ask again before concluding that
  no arbiter exists.
  \rationale{An arbiter that stays silent may be the passive one, or one waiting
  out a lease period under \reqref{R-0059}. A component that treats the silence as
  an absent arbiter and acts anyway creates the second leader the silence was
  preventing. A component therefore leads only on a grant, never on the absence
  of an answer.}
  \uncovered{no scenario has a component hear nothing from any voter and checks that it does not lead}
\end{requirement}

\begin{requirement}{R-0061}{Entitlement is held under a lease and lapses unless renewed}
  An instance's entitlement to act shall expire unless it is renewed, and an
  instance whose entitlement has expired shall stop acting at once.
  \rationale{A record that never expires states only that something was true at
  the moment it was written. Presence alone cannot be the test either: an
  instance that holds its connection open and stops renewing looks, from the
  outside, the same as a healthy one. An instance that keeps acting after its
  lease has run out is a second leader the moment a voter grants its peer a
  lease, so stopping is what makes the lease mean anything.}
  \uncovered{no scenario stops a leader renewing while it stays connected}
\end{requirement}

\begin{requirement}{R-0099}{Adjusting a machine's clock does not move an entitlement}
  The expiry of a lease shall be measured by a clock that only moves forward at a
  steady rate, and shall not be affected by the machine's notion of the time of
  day being corrected, stepped or set.
  \rationale{Time-of-day clocks are adjusted routinely, by a time service
  correcting drift, by a leap second, or by an operator. A lease measured against
  such a clock can expire while it is being renewed properly, which moves an
  entitlement for a reason unconnected to any component's health. An adjustment
  in the other direction is worse: a lease that should have expired goes on being
  honoured after its holder has gone.}
  \uncovered{no scenario adjusts a clock}
\end{requirement}

\begin{requirement}{R-0109}{An instance that is not making progress does not keep its entitlement}
  Holding an entitlement shall depend on the instance continuing to do the work
  the entitlement is for, and not only on its continuing to ask to keep it.
  \rationale{Renewing an entitlement requires no more than a running thread. An
  instance starved of processor time, or stalled in a way that leaves its timers
  running, renews on schedule and does nothing else. Neither presence nor renewal
  reports whether the instance is doing any work, so both report such an instance
  as healthy. The venue then stays with an instance that has stopped, and
  everything waits while nothing appears to be wrong.}
  \uncovered{no scenario starves an instance that keeps renewing}
\end{requirement}

\begin{requirement}{R-0115}{How long an entitlement lasts without renewal is a stated figure}
  The period after which an unrenewed entitlement lapses shall be a figure the
  venue states, and it shall be long enough that an instance doing its work is
  never mistaken for one that has stopped.
  \rationale{A period that is too short takes the entitlement away from a busy
  instance for being busy, which moves the venue's work to the other instance at
  the moment there is most of it. A period that is too long leaves the
  entitlement with an instance that has genuinely stopped, and nothing acts until
  it lapses. The figure alone shows neither error, so it has to be derived from
  how long the work takes.}
  \uncovered{no scenario tests the period, and none is stated to test}
\end{requirement}

\begin{requirement}{R-0062}{An instance already acting keeps its entitlement when another returns}
  Where an instance holds the entitlement and is still present, an instance
  rejoining the group shall be told to follow, whatever its configured name.
  \rationale{Preferring the configured primary at this point moves the
  entitlement for no reason but a name. It also moves the entitlement to the
  instance that has just failed, and away from one that is serving. The
  configured preference is a cold-start tie-break.}
  \covers{ha_test:24, ha_test:25}
\end{requirement}

\begin{requirement}{R-0063}{Losing one arbiter, or the witness, changes nothing}
  The failure of a single arbiter, or of the witness, shall not move any
  group's entitlement and shall not interrupt the venue.
  \rationale{Redundancy is what makes the arbiter available, so nothing needs to
  watch it. A design in which losing one arbiter mattered would need something to
  supervise the arbiter, and that supervisor would need supervising in turn.}
  \covers{ha_test:2, ha_test:4, ha_test:5, ha_test:33}
\end{requirement}

\begin{requirement}{R-0064}{A generation is never issued twice, and never goes backwards}
  Each grant of entitlement shall carry a generation greater than every
  generation the venue has used. Where every instance that could have issued one
  has restarted, the new generation shall still exceed every generation
  previously used, which requires that the number outlive the processes rather
  than any one of them.
  \rationale{The generation is what lets a receiver refuse a superseded sender. A
  generation held only in memory restarts at the beginning, so a pair that
  restarts together agrees a generation the venue has used before, and a message
  left over from the earlier one then compares as current.

  This is the requirement most exposed to a restart of everything at once. Where
  each instance keeps its own record, a group whose members all restart can still
  exceed every generation used, because each instance reads its own record and
  the highest wins. Where the record is kept in one place, the survival of that
  place is the whole of the guarantee. This chapter does not settle which of the
  two the venue uses. The requirement holds either way, and one of them has to be
  chosen deliberately.}
  \uncovered{no scenario restarts every instance of a group together}
\end{requirement}

\begin{gap}{}
  A generation can go backwards. An arbiter that has restarted, and has not been
  told the highest generation by the other arbiter because it has restarted too,
  can grant a new leader a generation below one already led in. Two instances
  still never act at once, but receivers that saw the higher generation refuse
  what the new leader sends. The new leader learns of the higher generation when
  its peer refuses it, and asks again above it. Where the peer is also down, only
  a receiver could tell it, and receivers do not yet do so. See
  \texttt{docs/availability/tla/findings.md}, section 11.5.
\end{gap}

\begin{requirement}{R-0065}{An instance whose entitlement has passed has what it sends refused}
  A receiver shall refuse a message carrying a generation older than the newest
  it has accepted for that group.
  \rationale{An instance whose lease has run out stops acting by itself
  (\reqref{R-0146}), before anything it sends could be refused. Refusing an older
  generation is the second defence, for anything that reaches a receiver all the
  same: a message delayed in transit, or a leader granted a generation below one a
  receiver has already seen. It needs no special hardware and no cooperation from
  the instance being refused.}
  \uncovered{a superseded instance stops before it sends, so no scenario makes one send anything for a receiver to refuse}
\end{requirement}

\subsection{The sequencer pair}
\label{sec:ha:enabled:sequencer}

Two sequencers run. One of them holds the entitlement to act. The other receives
everything the acting instance does, and waits.

\subsubsection{What state it holds}
\label{sec:ha:enabled:sequencer:state}

The pair holds the same state as the single sequencer in
\cref{sec:ha:disabled:sequencer}: the ordering, the log, the generation and each
session's numbering. The pair holds it twice, on two machines with two disks.

The follower builds its copy from the same records, in the same order, as the
leader writes them. A session's numbering reaches the follower by the same route
it reaches the leader, because gateways report to both instances. The follower
therefore holds what a restarted instance would have to be told.

\subsubsection{Where the copy that outlives it lives}
\label{sec:ha:enabled:sequencer:durable}

The copy lives on the other machine, which is the point of this configuration. A
failed disk, a stopped machine and a hung kernel each leave the venue's record of
what it has done intact, because the venue wrote a second copy as it wrote the
first.

\textbf{The rule that makes the second machine worth running governs when a
member is told.} An order becomes durable in two stages. The leader writes the
order, which permits the venue to give the order to the matching engine. The
follower then acknowledges the order, which permits the venue to give the
resulting execution report to the gateway. A member is therefore never told about
an order that exists on one machine only.

That ordering adds a round trip to the reporting path, and it is what gives the
venue the property the second machine is run for. Removing the round trip to
shorten the path would remove that property.

\subsubsection{How a restarted instance becomes current}
\label{sec:ha:enabled:sequencer:recovery}

A promoted follower is already current, or close to it. It holds the records and
has acknowledged them, and it needs only to apply what it has received and not
yet applied.

\textbf{A promotion recovers more than the restart in
\cref{sec:ha:disabled:sequencer} does.} A restarted sequencer recovers its
sequence number from its log, its generation from a file, and its sessions by
being told, which leaves it holding nothing for a session that was not bound at
the time. A promoted follower was told the same things as they happened, so it
holds the numbering of sessions that no gateway would report.

The pair also settles leadership between themselves when they can see each other,
rather than asking the arbiter. Both instances hold both instance identities and
both generations, so both compute the same answer from the same inputs. The
arbiter decides the case the pair cannot settle alone: one instance cannot see
the other, and a second claimant may exist that it cannot see.

\subsubsection{What the member is told}
\label{sec:ha:enabled:sequencer:member}

Nothing. The member sees a brief pause in responses and no break in numbering.

\begin{requirement}{R-0104}{A member is told nothing false while an entitlement is moving}
  Between the loss of the instance that was acting and the moment its replacement
  has caught up, a member may receive no answer to what it sends, and shall not
  receive an answer that is wrong. An order shall not be rejected as unknown, and
  an order shall not be reported accepted, on the basis of state that is not yet
  current.
  \rationale{The interval is short, but it is not zero, and a member cannot tell that it is in
  one. Silence during the interval is a delay, and a member is built to tolerate a
  delay. An answer given from state that is still being assembled is a statement
  the venue cannot support, and the member acts on that statement.}
  \uncovered{no scenario sends orders during a failover and checks the answers}
\end{requirement}

\begin{requirement}{R-0066}{A member is never told of an order held on one machine only}
  An execution report shall not be sent to a member until the record of the order
  it concerns has been acknowledged by a second instance.
  \rationale{An order that exists in one place is lost when that place is lost, so telling a
  member about such an order commits the venue to something it may not be able to
  honour. This is the property the second machine is run for, and the other
  requirements in this section depend on it.}
  \uncovered{no scenario asserts that a report waits for the follower}
\end{requirement}

\begin{requirement}{R-0067}{Losing a machine does not lose a committed order}
  An order committed before a machine was lost shall be processed and answered by
  the surviving instance.
  \rationale{Committing the order is the venue taking it. Losing one of the two machines
  holding that order is the failure the pair exists to absorb.}
  \uncovered{no scenario asserts an order committed before the kill is answered after it}
\end{requirement}

\begin{requirement}{R-0068}{The venue keeps trading when the acting sequencer is lost}
  Members shall continue to be able to place orders across the loss of the
  sequencer instance that was acting, without reconnecting and without being
  told to retry.
  \rationale{This is the most visible reason for running two machines, and a member would
  notice immediately if it did not hold. It is stated separately from the
  requirements about what survives because continuing to trade and keeping what was
  taken are different properties. A venue can meet this requirement while failing
  every other requirement in this section.}
  \covers{ha_test:1}
\end{requirement}

\begin{requirement}{R-0069}{A member's session survives a failover}
  A member's numbering, and the orders open under its session, shall be
  unaffected by leadership moving between sequencers.
  \rationale{A failover the member can detect has not achieved what it was for. The venue
  tells the follower each session's position as that position advances, for this
  reason.}
  \uncovered{no scenario fails the sequencer over under a live session}
\end{requirement}

\begin{requirement}{R-0070}{A follower that falls behind stops the venue confirming, rather than being dropped}
  Where the second instance cannot keep up, the venue shall stop sending
  execution reports rather than continue without it.
  \rationale{Dropping the follower to keep the reporting path moving abandons \reqref{R-0066}
  without reporting it, at the moment that requirement matters most. Stopping the
  reports instead leaves members seeing orders received and no reports, which a
  member can see and recover from. A member cannot see the venue running on one
  machine.}
  \uncovered{no scenario makes the follower lag}
\end{requirement}

\begin{requirement}{R-0100}{Waiting for a disk or for a peer does not stop a component working}
  A component shall continue to process what arrives while it is waiting for a
  write to complete or for a second instance to acknowledge one, and shall
  continue to renew its entitlement and to answer enquiries about its health
  throughout.
  \rationale{The requirements in this section make a component wait for a disk and for a peer,
  and both can be slow for reasons unconnected to the venue. A component that stops
  processing while it waits also stops answering members, stops replying to its
  peer and stops renewing its lease. From outside, such a component looks the same
  as one that has stopped, so the venue restarts it, and a slow disk becomes a
  restart rather than a delay.}
  \uncovered{no scenario makes a disk or a peer slow}
\end{requirement}

\begin{requirement}{R-0110}{How far behind a second instance may fall is a stated figure}
  The point at which a second instance is too far behind shall be a figure the
  venue states, expressed in what has not been acknowledged rather than in
  elapsed time alone.
  \rationale{\reqref{R-0070} requires the venue to stop confirming rather than proceed without
  the second instance. Neither implementing nor testing that requirement is possible
  without knowing when the point is reached. A figure expressed only as a duration
  moves with the rate orders arrive, so a venue under load reaches it after far
  fewer records than a quiet venue does.}
  \uncovered{no scenario makes a second instance lag, and no figure says when it is too far}
\end{requirement}

\begin{requirement}{R-0071}{A sequencer that finds its own record damaged stops rather than serving from it}
  An instance that cannot read its log whole shall not act on the part it can
  read, and the instance holding an undamaged copy shall act instead.
  \rationale{This is the answer that \cref{sec:ha:disabled:recovery} states the disabled
  configuration does not have. An instance serving from a partial log assigns
  numbers it has assigned before, and answers orders other than the ones it
  committed.}
  \uncovered{no scenario damages a log}
\end{requirement}

\begin{requirement}{R-0072}{A pair that can see each other settles leadership between themselves}
  Where both instances are present and neither holds the entitlement, they shall
  reach the same answer without referring it, and shall reach it whichever
  started first.
  \rationale{Both instances hold both identities and both generations, so both compute the
  same result from the same inputs. Asking the arbiter instead makes the venue wait
  on a component that has just started itself, and that holds less information than
  either instance.}
  \covers{ha_test:36, ha_test:37}
\end{requirement}

\subsection{The matching engine pair}
\label{sec:ha:enabled:me}

Two matching engines run. One of them acts. The other holds a replica of the set
of open orders, built from the same records in the same order.

\subsubsection{What state it holds}
\label{sec:ha:enabled:me:state}

The pair holds the set of open orders and the sequence number that set is current
to, as \cref{sec:ha:disabled:me:state} describes, and holds them twice. The
engine keeps no log of its own in either configuration. The sequencer holds the
record of what was taken, and what the engine holds follows from that record.

\subsubsection{Where the copy that outlives it lives}
\label{sec:ha:enabled:me:durable}

The copy lives in the other instance, kept current by the same stream.
\Cref{sec:ha:disabled:me:durable} states that this copy does not exist when high
availability is disabled, which is why that configuration needs a checkpoint.

\subsubsection{How a restarted instance becomes current}
\label{sec:ha:enabled:me:recovery}

The promoted instance reports the sequence number its replica has reached,
receives everything committed after that number, applies it silently, and begins
acting.

\textbf{This is the same operation as \cref{sec:ha:disabled:me:recovery}
describes.} A restarting instance learns its starting number from a checkpoint,
and a promoted instance takes it from the replica it has been maintaining.
Nothing else about the operation differs, so the mechanism built for the disabled
configuration is the mechanism this configuration uses.

The engine differs from the sequencer in one respect. A sequencer pair settles
leadership between themselves in both directions, because both instances hold
both generations at the same moment. A matching engine only ever \emph{gives}
leadership up. It adopts the follower position on its peer's assertion, takes the
generation it is given, and issues no generation of its own. An instance that
gives up leadership cannot produce a second acting instance, whatever its peer
has got wrong, which is why acting on the peer's assertion alone is safe here.

\subsubsection{What the member is told}
\label{sec:ha:enabled:me:member}

\begin{requirement}{R-0073}{An open order survives a failover}
  An order open before leadership moved shall be open after it, on the same
  terms, and shall remain cancellable by the member that placed it. Where the set
  cannot be recovered, \reqref{R-0020} applies instead and every order the venue
  can still name is cancelled to its member.
  \rationale{\reqref{R-0018} requires this of a restart. The configuration with high
  availability enabled exists to give a member a better outcome, so a member must
  not keep its orders when a process dies and lose them when a machine dies.

  The two requirements form an order of preference rather than a choice. The venue
  recovers what it can recover, and reports what it cannot. The member is told
  exactly where it stands in either case, and the venue cancels only where it does
  not have the orders.}
  \uncovered{scenario 21 asserts the fallback, and nothing asserts this}
\end{requirement}

\begin{gap}{}
  \textbf{The venue does not meet this requirement as it is configured, and says
  so.} What a promoted matching engine does with the orders it inherits is a
  stated policy, \texttt{order\_book.open\_orders\_on\_promotion}, and it is set
  to \texttt{cancel} in every environment. Every promotion therefore takes the
  fallback path: each order is cancelled, its member is told, and the venue leads
  an empty book.

  The engine can give the other answer, because \texttt{keep} carries the book
  across. The mechanism exists, and what is missing is the grounds for trusting it.
  A promoted instance assembles its book from two sources, and \textbf{one of them
  is checked and the other is not}. The catch-up applied before the instance acts
  is accounted for, which is \reqref{R-0101}. The replica the instance maintained
  by \texttt{BookUpdate} while it was the follower is not accounted for: it applies
  those updates as they arrive, and nothing establishes that none was missed.
  Keeping a book that is short of an order is worse than cancelling every order,
  because the member is told nothing and believes it holds an order that the venue
  does not.

  The setting therefore stands at the outcome that is worse and certain. Changing
  it requires the same check to be applied to the replication stream that
  \reqref{R-0101} applies to the catch-up.
\end{gap}

\begin{requirement}{R-0074}{An instance that gives up leadership cannot thereby create a second one}
  An instance shall adopt the follower position only on its peer's assertion that
  the peer leads, only while itself holding no position, and only where the
  generation offered is not older than one it has already seen.
  \rationale{An instance that gives up a position cannot produce two acting instances,
  whatever its peer has got wrong. A one-way announcement is therefore sufficient
  here, where the sequencer needs an exchange.}
  \uncovered{no scenario offers a stale generation to an instance holding no position}
\end{requirement}

\begin{requirement}{R-0075}{The venue keeps working when the acting engine is lost}
  Members shall continue to be able to place and cancel orders across the loss
  of the matching engine instance that was acting.
  \rationale{Cancelling matters here as much as placing. A member that can place orders but
  cannot cancel the orders it already holds is worse off than a member that can do
  neither, because it accumulates exposure that it has no way to reduce.}
  \covers{ha_test:16}
\end{requirement}

\begin{requirement}{R-0076}{Orders taken while no engine was acting are applied on promotion}
  An order committed while no instance was acting shall be applied and answered
  once one is.
  \rationale{The venue took those orders, so it must answer each of them. The disabled
  configuration cannot serve this case, because it has no second instance holding a
  position to catch up from.}
  \uncovered{no scenario asserts that deferred orders are answered after a promotion}
\end{requirement}

\subsection{The publisher and recorder pairs}
\label{sec:ha:enabled:publishing}

Two publishers run, and two recorders.

\subsubsection{What state it holds}
\label{sec:ha:enabled:publishing:state}

The publisher holds what it has been given, and holds no position on behalf of
any subscriber. A subscriber presents its own position when it connects, as
\cref{sec:ha:disabled:mep:state} describes. The recorder holds one number: how
far it has published and had confirmed.

\subsubsection{Where the copy that outlives it lives}
\label{sec:ha:enabled:publishing:durable}

Both instances of a publisher receive the same records, identified by the same
sequence numbers. Those numbers are the sequencer's, and the publisher assigns
none of its own. \textbf{A subscriber's position therefore means the same thing
at either instance}, so a subscriber moving to the surviving instance needs no
translation, and the two publishers need no handover between them.

The recorder's position needs less still. The recorder writes it only after the
bus confirms the publish, so the position is a lower bound on what has been
published. A promoted instance resuming from a stale copy publishes some records
again and loses none, and the consumers downstream are required to handle the
duplicates.

\subsubsection{How a restarted instance becomes current}
\label{sec:ha:enabled:publishing:recovery}

Other components tell both of them. Each subscriber tells a publisher where that
subscriber had reached. A recorder reads its own position and presents it to a
publisher. Neither needs a handover from the instance it replaced, so neither
pair needs the coupling that the sequencer pair has.

\subsubsection{What the member is told}
\label{sec:ha:enabled:publishing:member}

The publisher tells a member nothing. A failure of the recorder reaches members
only through \reqref{R-0053}, and running two recorders is what keeps it from
reaching them at all.

\begin{requirement}{R-0077}{A subscriber resumes at the surviving instance without loss}
  Where the instance a subscriber was reading from is lost, the venue shall
  supply every record after the position that subscriber presents, from the
  instance it moves to.
  \rationale{A subscriber's position is expressed in the venue's own sequence numbers rather
  than in numbers a publisher assigns, so either instance can act on it. The two
  publishers hand nothing over.}
  \uncovered{no scenario moves a subscriber between publisher instances}
\end{requirement}

\begin{requirement}{R-0078}{A promoted recorder republishes rather than skips}
  Where publishing to the bus is taken over by another instance, the venue shall
  resume from a position no later than the last publish the bus confirmed, and
  shall publish again whatever it cannot establish was confirmed.
  \rationale{The position is a lower bound, so resuming from a stale position costs duplicate
  records, and resuming from a position later than the confirmed one costs records
  entirely. \reqref{R-0049} permits the first and forbids the second.}
  \uncovered{no scenario, and none is possible until the recorder exists}
\end{requirement}

\begin{requirement}{R-0079}{Losing one publisher or one recorder does not halt the venue}
  The loss of a single publisher instance, or of a single recorder instance,
  shall not stop members trading.
  \rationale{\reqref{R-0053} halts the venue when no recorder is publishing. Running two
  recorders is what makes that condition rare. Halting on the loss of one of them
  would produce the outage that the second recorder exists to prevent.}
  \uncovered{no scenario, and none is possible until the recorder exists}
\end{requirement}

\subsection{A process dies while its peer is present}
\label{sec:ha:enabled:process_death}

A component dies on a machine that is otherwise healthy, and the launcher
restarts it, while a second instance of the same component runs elsewhere. Both
configurations apply to this case.

\subsubsection{What is at risk}
\label{sec:ha:enabled:process_death:risk}

The venue's ability to trade is not at risk, because the peer can act and trading
continues whatever the venue decides. What is at risk is \textbf{the pair}.

If the peer promotes, one instance serves the venue until somebody restores the
other, and the venue is exposed to a second failure for as long as that takes. If
the launcher restarts the dead instance instead, the pair is whole again in
seconds and the venue loses nothing.

Nothing the venue does can restore a machine that has died, so promoting the peer
is the only remedy for that failure. The venue can restore a process that has
died on a living machine, so promoting the peer in that case gives up a recovery
that was available. The venue must therefore treat the two failures differently.

\subsubsection{What the venue must do}
\label{sec:ha:enabled:process_death:behaviour}

\begin{requirement}{R-0080}{A process death on a healthy machine does not cost the pair}
  Where a component dies and is restarted on a machine that has not failed, the
  venue shall end with two instances running, as it began.
  \rationale{The measure of this requirement is whether the pair is whole again, rather than
  whether the venue kept trading. A promotion keeps the venue trading on its own,
  and leaves one instance down with nothing watching for a second failure.}
  \uncovered{no scenario asserts the pair is whole afterwards, only that trading continued}
\end{requirement}

\begin{requirement}{R-0081}{The peer does not promote while a restart could still succeed}
  An instance that has lost contact with the acting instance shall wait, before
  seeking to act, for at least as long as a supervised restart takes.
  \rationale{A shorter wait fails the venue over for a failure that did not require it. A much
  longer wait tolerates a genuinely dead machine for longer than it should. The
  period is therefore the time a restart takes plus a margin. That time is not
  known until restarts have been timed, so the period is measured rather than
  chosen, and it must be measured again whenever the components change.
  The peer cannot take over until its own promise to the acting instance and the
  arbiter's have both run out, which is at least one lease period after the last
  renewal. A restarted instance that kept its promises on disk asks again as soon
  as it starts, so a restart that completes within the lease period keeps the
  lead.}
  \covers{ha_test:26, ha_test:27, ha_test:28}
\end{requirement}

\begin{requirement}{R-0082}{An instance that keeps failing does not disturb a healthy peer}
  A restarted instance shall not take the entitlement to act from a peer that
  holds it and is present.
  \rationale{Without this requirement, an instance that fails repeatedly takes leadership on
  each restart, dies, and hands leadership back, interrupting a peer that was
  working correctly once per failure. \reqref{R-0062} states the arbiter's side of
  this rule, and this requirement states the instance's side.}
  \covers{ha_test:29, ha_test:30, ha_test:31}
\end{requirement}

\begin{requirement}{R-0125}{An instance that interrupts the venue twice is not restarted again}
  Where restarting an instance has already cost the venue one interruption, a
  second interruption from that instance shall leave it stopped rather than
  restarted.
  \rationale{An interruption is a period in which a member can neither place an order nor
  withdraw one, and is told nothing while it lasts. The venue can account for one
  interruption: a component failed, and recovery took as long as it took. A second
  interruption from the same instance shows that the instance is unfit to serve,
  and restarting it again produces a third.

  Stopping the instance is what allows the pair to do its work. An instance that
  the supervisor keeps restarting returns before the arbiter can grant the
  entitlement to its peer, so the peer never acts. The venue then takes an
  unbounded series of interruptions instead of the single failover the pair was
  built for. Leaving the failing instance stopped produces the outcome the member
  should have had at the first fault.

  \textbf{The count is taken over a session rather than over a period.} A time
  window cannot be explained to the member it affects, because it makes two
  failures count together when they fall close in time and separately when they do
  not. The requirement is stated of interruptions rather than of failed starts for a
  related reason: an instance that dies every few minutes starts correctly every
  time.}
  \uncovered{no scenario flaps an instance and asserts it is left stopped}
\end{requirement}

\begin{requirement}{R-0126}{The venue does not run with no instance of a component without saying so}
  Where the last instance of a component is stopped and not restarted, the venue
  shall report that it cannot perform that component's work, rather than continuing
  in a state where nothing does it.
  \rationale{Both instances may share the cause of the failure: an input that kills whatever
  reads it, a build, or a configuration both instances load. The peer then meets
  that cause in its turn and is stopped in its turn. The venue has no instance of
  that component left, and nothing is trying to produce one.

  This failure looks like health. Every process that remains is running normally,
  no alert names a component, and only somebody who counts the processes can see
  what is missing. Reporting the condition costs an announcement. Not reporting it
  costs the time between the last instance stopping and somebody noticing, and that
  time is unbounded.}
  \uncovered{no scenario stops every instance of a component and asserts the venue says so}
\end{requirement}

\subsubsection{When the restart and the promotion race}
\label{sec:ha:enabled:process_death:race}

The two remedies run at the same time, and neither detects the other. The
launcher restarts the instance that died while the peer counts down to asking
whether it may act. \textbf{Exactly one instance ends up acting, whichever of the
two happens first.}

\textbf{The restart completes first.} The returning instance holds no position
and asks which position it holds. No instance holds the entitlement, so the
arbiter grants it to the returning instance, which resumes acting. The peer is
still waiting, finds the entitlement held, and remains the follower. The
entitlement did not move, and the pair is whole.

\textbf{The peer asks first.} The arbiter grants the peer the entitlement, and
the peer begins acting. The returning instance asks, finds the entitlement held
by an instance that is present, and becomes the follower. Neither its configured
name nor its having held the entitlement minutes earlier affects that answer.

\textbf{Both ask at once.} One party answers both claims. The arbiter grants the
entitlement to one instance, records what it granted, and refuses the other claim
against that record.

In every ordering, exactly one instance ends up acting, and no instance has to
detect what the other was doing. That is what makes the outcome certain rather
than merely likely.

\begin{requirement}{R-0090}{A returning instance assumes nothing and asks}
  An instance that has restarted shall hold no position until it is told which
  it holds, whatever position it held before and whatever it is configured as.
  \rationale{An instance that assumes it leads, because it led before or because of its
  configured name, becomes a second acting instance whenever the peer has already
  been promoted. Nothing has to go wrong with the peer for that to happen.}
  \covers{ha_test:24, ha_test:25}
\end{requirement}

\begin{requirement}{R-0091}{A returning instance that finds the position vacant takes it}
  Where an instance returns and no instance holds the entitlement, it shall take
  it rather than wait.
  \rationale{A returning instance holds no position and asks, which is what stops a restart
  producing a second acting instance. This requirement states the other half of
  that rule. Where no instance holds the entitlement, the returning instance must
  take it. Without this, a restart that completed before its peer's promotion would
  leave both instances following and none acting.}
  \covers{ha_test:26, ha_test:27, ha_test:28}
\end{requirement}

\begin{requirement}{R-0092}{Simultaneous claims are resolved by one party, not by the claimants}
  Where a returning instance and its peer both seek the entitlement at once, at
  most one shall be granted it, and neither shall be required to have detected
  the other.
  \rationale{Two instances that must detect each other in order to stay out of each other's
  way will eventually fail to do so, and they will fail at the moment they are
  least able to communicate. Referring both claims to one party removes the need to
  detect anything.}
  \uncovered{no scenario times a restart to collide with a promotion}
\end{requirement}

\subsection{Losing a machine}
\label{sec:ha:enabled:machine_loss}

A machine stops. This is the failure the configuration exists for, and it is
always a combined failure, because a machine hosts several components and losing
it loses all of them at the same instant.

\subsubsection{What is at risk}
\label{sec:ha:enabled:machine_loss:risk}

Every component that machine was running is at risk at once, together with every
position any of them held. Nothing on that machine can be restarted, so for each
of those groups the only remedy is the instance on the other machine.

Two requirements follow that do not arise when a single process dies.
\textbf{The surviving machine must be able to carry the whole venue}, rather than
the share it was carrying. \textbf{The components that decide which instance may
act must not all have been on the machine that stopped}, or the venue loses the
ability to decide at the moment it most needs to.

\subsubsection{What the venue must do}
\label{sec:ha:enabled:machine_loss:behaviour}

\begin{requirement}{R-0083}{The venue survives the loss of any one machine}
  Members shall continue to trade across the loss of any single machine,
  whichever machine it is and whatever it was running.
  \rationale{This requirement is stated of any machine rather than of a nominated one. A
  deployment that survives losing the secondary and not the primary does not meet
  it, and nothing distinguishes the two deployments until the wrong machine fails.}
  \covers{ha_test:48}
\end{requirement}

\begin{requirement}{R-0084}{No single machine holds enough of the venue to prevent a decision}
  The venue shall be deployed so that the loss of any one machine leaves it
  still able to decide which instance may act.
  \rationale{An arbiter pair and its witness deployed on one machine share that machine's
  failure. Losing the machine then leaves the venue with no arbitration at all.}
  \uncovered{no scenario deploys instances badly, and nothing would report it}
\end{requirement}

\begin{requirement}{R-0111}{A deployment that could not survive losing a machine is reported when it starts}
  Where instances are placed such that the loss of one machine would leave the
  venue unable to decide which instance may act, that shall be reported when the
  venue starts rather than discovered when the machine is lost.
  \rationale{\reqref{R-0084} is a property of where the processes were placed. A deployment
  that breaks it behaves the same as one that does not, until the wrong machine
  fails. The venue can check the placement when the instances announce themselves,
  and that is the only moment before the failure at which anybody would look.}
  \uncovered{no scenario deploys instances badly, and nothing would report it}
\end{requirement}

\begin{requirement}{R-0085}{The surviving machine can carry the whole venue}
  A machine shall be able to run everything the venue requires of it after its
  partner is lost, and not merely its own share.
  \rationale{Where each machine is sized to carry only its own share, the surviving machine
  takes on a load it cannot carry and fails in turn. The venue then delays the
  outage rather than avoiding it, and loses the second machine as well as the
  first.}
  \uncovered{no scenario measures the survivor under the whole load}
\end{requirement}

\subsection{A partition}
\label{sec:ha:enabled:partition}

Both instances are running, and neither can see the other. Both are healthy, and
neither can establish whether the other has died or has merely become
unreachable.

This is the condition the arbiter exists for. It is the only condition in this
chapter where both available actions, promoting and declining, produce a wrong
outcome in one of the two cases.

\subsubsection{What is at risk}
\label{sec:ha:enabled:partition:risk}

An instance that cannot see its peer has two explanations available and no way to
choose between them. If it promotes itself and the peer is alive, both instances
act and their states diverge, and nothing afterwards can establish which of them
was right. If it declines and the peer has died, no instance acts. That second
failure is harder to notice, because no process crashes and nothing logs an
error, so somebody finds it by asking why trading has stopped.

\subsubsection{What the venue must do}
\label{sec:ha:enabled:partition:behaviour}

\begin{requirement}{R-0086}{An instance that cannot see its peer refers the question rather than deciding it}
  Where an instance cannot reach its peer, it shall ask the arbiter which
  instance may act, and shall not decide for itself.
  \rationale{This is the one case an instance cannot settle alone, because it is the only case
  in which a second claimant may exist that the instance cannot see. Where an
  instance can see its peer, the two settle the question between themselves under
  \reqref{R-0072}. The two rules cannot both apply at once, so no group has two
  parties entitled to decide.}
  \uncovered{no scenario partitions a pair}
\end{requirement}

\begin{requirement}{R-0087}{A partition does not produce two instances acting}
  A partition of any duration shall not result in both instances of a group
  acting.
  \rationale{\reqref{R-0058} states this property generally. This subsection restates it
  because a partition is the condition that produces the failure, and because a
  design can satisfy the general statement in every case a test produces and still
  fail here.}
  \uncovered{no scenario partitions a pair}
\end{requirement}

\subsubsection{When only one side can see the other}
\label{sec:ha:enabled:partition:asymmetric}

A partition need not be symmetric. One instance may be able to reach its peer
while the peer cannot reach it, and the two rules for settling an entitlement
then both apply at once: the instance that can see its peer settles it locally
under \reqref{R-0072}, while the instance that cannot refers it under
\reqref{R-0086}.

The two rules cannot both apply only while visibility is the same in both
directions, and here it is not. What settles this case is that an instance taking
an entitlement reports it immediately, rather than at its next scheduled report.
The arbiter therefore already holds a record of a present holder, and confirms
that holder rather than granting the entitlement again.

The window is therefore as long as a new holder takes to report itself, rather
than as long as the partition. The window is narrowed rather than removed, which
is a weaker guarantee than the venue has where one party alone issues
entitlements. A scenario should produce this case deliberately.

\subsubsection{What this costs, stated plainly}
\label{sec:ha:enabled:partition:cost}

The venue does not stop an instance on the losing side of a partition. That
instance keeps running, goes on treating itself as the leader, and is refused at
every interaction it attempts, under \reqref{R-0065}. It corrupts nothing,
because no receiver accepts what it sends. It also does not stop, and it may hold
resources until somebody intervenes.

The venue accepts that cost. Removing the instance instead would require
authority over the machine it runs on: a separate management network, hardware
power control, and credentials to use them. Those are properties of a deployment
rather than of this software, and a developer's machine has none of them, so such
a mechanism could only ever be exercised in one environment. Refusing everything
a superseded instance sends requires none of them, and is exercised wherever the
venue runs.

\subsection{Failback}
\label{sec:ha:enabled:failback}

An instance that failed is repaired and returns. This subsection settles whether
the entitlement to act returns with it.

\textbf{The entitlement stays where it is.} The instance that has just returned
is the one that recently failed, and the instance holding the entitlement is
serving the venue. Moving the position from the second to the first, on the
strength of a name fixed in configuration when the deployment was built, changes
which instance acts for a reason unconnected to either instance's fitness.

An operator may still move leadership deliberately. A machine may need to be
taken out of service, or a deployment rebalanced, and moving the position on
purpose is an ordinary operation. The venue must not move it automatically.

\begin{requirement}{R-0088}{The entitlement does not return to an instance merely because it has come back}
  A repaired instance rejoining a group shall take the follower position while
  the other instance holds the entitlement and is present, whatever its
  configured name.
  \rationale{Configured names order a cold start and decide nothing else. An automatic return
  moves the position away from an instance that is working, to the instance that
  recently failed, at a moment nobody chose.}
  \covers{ha_test:24, ha_test:32}
\end{requirement}

\begin{requirement}{R-0089}{Leadership can be moved deliberately, and moving it loses nothing}
  An operator shall be able to move the entitlement between instances of a group,
  and doing so shall lose no order, no report and no session.
  \rationale{Without this requirement, taking a machine out of service means provoking the
  failure the configuration exists to survive. An operator who must kill a healthy
  process to move a position will eventually kill the wrong one.}
  \uncovered{no scenario moves an entitlement on purpose, and nothing provides for it}
\end{requirement}

\subsection{Running degraded}
\label{sec:ha:enabled:degraded}

Every failure in this section leaves the venue in some state. For several of
them, the venue is working and can no longer survive the next failure.

A group whose second instance has not yet been restarted serves from one
instance. A venue that has lost both arbiters goes on trading, and cannot move an
entitlement if it needs to. Neither condition stops anything, neither produces an
error, and an observer outside the venue cannot distinguish either from a venue
in full health.

\textbf{The venue must go on running in these states.} Stopping when it loses
redundancy would turn the loss of a second instance into the outage that second
instance existed to prevent. The venue must also report that it is running in
such a state, because the exposure lasts until somebody restores what was lost.

\subsubsection{What the venue must do}
\label{sec:ha:enabled:degraded:behaviour}

\begin{requirement}{R-0093}{The venue trades while a group is reduced to one instance}
  The loss of one instance of a group shall not stop members trading, and shall
  not stop the surviving instance acting.
  \rationale{A pair reduced to one instance has lost its protection against the next failure,
  and has lost nothing else. Refusing to work in that state produces an outage now
  instead of risking one later.}
  \covers{ha_test:3, ha_test:6}
\end{requirement}

\begin{requirement}{R-0094}{A group running without its second instance is knowable}
  The venue shall make it possible to establish, at any moment, which groups are
  running with a second instance and which are not.
  \rationale{The period without resilience lasts from the failure until something restores the
  missing instance, and its length is the venue's exposure to a second failure.
  Automatic restart makes that period short, and also makes it invisible. Somebody
  downstream eventually notices a component that dies and stays dead, and nobody
  notices one that is restarted repeatedly. The venue must therefore report what it
  recovers automatically.}
  \uncovered{no scenario asks whether a group is running single, and nothing answers}
\end{requirement}

\begin{requirement}{R-0095}{An instance may act with no arbiter reachable, and says that it did}
  Where no arbiter can be reached and an entitlement must be settled, an instance
  may take it by the rule that orders a cold start, and shall record that it did
  so without arbitration.
  \rationale{Refusing to act because the arbiter is unreachable stops the venue for the
  failure of a component that holds no venue state and sits off the order path. It
  is also the one moment at which the venue has no protection against a second
  claimant, so the venue must record that it acted without arbitration.}
  \covers{ha_test:8, ha_test:9, ha_test:35, ha_test:39}
\end{requirement}

\subsection{How the design is checked}
\label{sec:ha:enabled:model_checking}

The scenarios in \texttt{scripts/ha\_test.py} run the venue and inject failures.
Each scenario tries one sequence of events. The failures that matter most in this
section depend on the order in which events happen: a message arriving just
after a restart, or two timeouts firing close together. A scenario that passes
says nothing about the orders it did not try.

\textbf{The design is therefore also checked by model checking.} The design is
written down in TLA+, a language for describing a system precisely enough for a
program to check it. The description states what each component holds, which
steps can change it, and which properties must always hold. The model checker,
TLC, then visits every state that a small version of the system can reach, in
every order of events, and reports any sequence of steps that breaks a property.
That sequence is called a counterexample.

The check has two limits. It covers a small version of the system: two instances,
one arbiter or two, and one or two failures at a time. It checks the design as
written in TLA+, so each property it confirms holds in the code only as far as
the TLA+ description matches the code.

\subsubsection{What is checked}
\label{sec:ha:enabled:model_checking:scope}

Two TLA+ descriptions model the design the code implements:

\begin{itemize}
  \item the two sequencer instances and the arbiter they report to; and
  \item the two arbiters and the witness, which decide which arbiter is active.
\end{itemize}

Both include crashes and restarts of every party, the failure of any link, and
messages that arrive after the state they describe has changed.

Model checking found eight defects in that design. The code now addresses six of
them in full and part of one more. The properties that now hold within the limits
checked are these:

\begin{itemize}
  \item no two leaders ever lead in the same epoch, so a receiver can always tell
        them apart;
  \item two leaders that can hear each other always resolve to one; and
  \item two active arbiters that can hear each other always resolve to one.
\end{itemize}

\textbf{Two instances can still both lead while a partition lasts.} An instance
that can reach no arbiter promotes itself, as \reqref{R-0095} permits. When the
real cause is that the instance has been cut off, the pair has two leaders until
the partition heals. The arbiter pool can likewise have two active arbiters during
a partition, because it does not yet decide by majority.

A replacement design removes promotion without an arbiter. In the replacement, an
instance leads only while a majority of three voters has granted it a lease that
has not run out. It is a proposal, described in
\texttt{docs/availability/majority\_leases.md}, and nothing in the code implements
it. A third TLA+ description specifies it. Model checking that description found
no state in which two instances act as leader at once. It also showed that each
rule of the replacement is needed: with any one rule removed, TLC finds two
instances acting as leader at once within a few steps.

The descriptions, the counterexamples and the full results are in
\texttt{docs/availability/tla/}, with the results in \texttt{findings.md}.

\subsubsection{What the venue's build must do}
\label{sec:ha:enabled:model_checking:behaviour}

\begin{requirement}{R-0145}{The build reproduces every recorded counterexample}
  The build shall rerun each counterexample recorded against a TLA+ description of
  the high availability design, and shall fail if any of them no longer breaks the
  property it is recorded against.
  \rationale{A counterexample is the evidence that a rule of the design is needed.
  A later change to the TLA+ description can stop a counterexample failing, and
  the description then no longer shows what the recorded results say it shows.
  Nothing else in the build would notice. Rerunning every counterexample takes a
  few seconds.}
  \covers{tla:lease-1-holder-counts-from-arrival, tla:lease-2-candidate-ignores-own-grant,
  tla:lease-3-restart-does-not-wait, tla:lease-4-degraded-promotion,
  tla:lease-5-epoch-regresses, tla:lease-6-no-yield-livelock, tla:lease-7-no-peer-vote}
\end{requirement}

The counterexamples reproduced today are those of the replacement design.
\texttt{scripts/tla\_trace\_pages.py} reruns them at install time, and writes a web
page that steps through each one.


